% !TeX spellcheck = de_DE
% !LuaLaTeX
\documentclass[12pt,a4paper,listof=totocnumbered,parskip=half]{scrreprt}
\usepackage{scrhack}
\usepackage[ngerman]{babel}
\usepackage[T1]{fontenc}
\usepackage[utf8]{inputenc}
\usepackage{lmodern}
\usepackage{microtype}
\usepackage{geometry}
\geometry{a4paper, top=27mm, left=20mm, right=20mm, bottom=35mm, headsep=10mm, footskip=12mm}
\usepackage{setspace}
\usepackage{amsmath,amssymb,amsthm,mathrsfs,amsfonts}

\usepackage{csquotes}
\usepackage{booktabs}

\usepackage{graphicx}
\graphicspath{ {./img/} }
\usepackage{wrapfig}
% \usepackage{layouts}
\usepackage{textcomp} 
\usepackage[pdftex,dvipsnames]{xcolor}  % Coloured text etc.
\usepackage{svg}
\svgsetup{inkscapelatex=false, inkscapearea=drawing}
\usepackage[layout={margin,index},draft]{fixme}
\usepackage{rotating}
\usepackage{siunitx}

\usepackage[backend=biber,minbibnames=2,maxbibnames=2,maxcitenames=1,mincitenames=1,style=alphabetic]{biblatex}
\setcounter{biburlnumpenalty}{9000}
\setcounter{biburlucpenalty}{9000}
\setcounter{biburllcpenalty}{9000}

% new stretchable space between characters
\setlength{\biburlnumskip}{0mu plus 1mu}
\setlength{\biburlucskip}{0mu plus 1mu}
\setlength{\biburllcskip}{0mu plus 1mu}
\renewcommand{\namelabeldelim}{\addnbspace}

\addbibresource{Recherche/My Collection.bib}
\defbibheading{Literatur}{\section{Literaturverzeichnis}} 
\defbibheading{Quellen}{\subsection*{Quellenverzeichnis}} 

\DeclareSourcemap{
  \maps{
    \map{
      \step[fieldsource=url,
            match=\regexp{\$\\sim\$},
            replace=\regexp{\~}]
    }
  }
}

\usepackage{tikz}
\usetikzlibrary{calc,positioning, shapes, petri, automata,spy}
\usetikzlibrary{external}
\usepackage{pdfpages}
\usepackage{subcaption}
\usepackage{standalone}
\usepackage{multirow,tabularx}
\usepackage{xifthen}
\usepackage[hidelinks]{hyperref}
\usepackage[acronym,shortcuts,toc]{glossaries-extra}
%\usepackage{glossary-longbooktabs}
\setabbreviationstyle[acronym]{long-short} 

\renewcommand{\glsxtrregularfont}[1]{\textcolor{violet}{\glsxtrifwasfirstuse{\emph{#1}}{#1}}}
\renewcommand{\glsxtrabbreviationfont}[1]{\textcolor{teal}{\glsxtrifwasfirstuse{\emph{#1}}{#1}}}
% or
%\renewcommand{\glsxtrregularfont}[1]{\textcolor{violet}{\ifglsused{\glslabel}{#1}{\emph{#1}}}}
%\renewcommand{\glsxtrabbreviationfont}[1]{\textcolor{teal}{\ifglsused{\glslabel}{#1}{\emph{#1}}}}

\renewcommand*{\glsxtrsetupfulldefs}{%
  \renewcommand*{\glsxtrifwasfirstuse}[2]{##2}%
}

\usepackage{mathtools}


\usepackage{pgfplots}
\pgfplotsset{
	grid=major,
	width=\textwidth,
	height=5cm,
	compat=1.18,
	no markers,
  table/search path={performance/measurements},
}
\usepgfplotslibrary{external}
\tikzexternalize[prefix=tikz/,optimize command away=\includepdf]
%\tikzset{external/force remake}

\newcommand{\unsim}{\mathord{\sim}}
\newcommand{\e}{\ensuremath{\mathrm{e}}}

\newcommand{\R}{\ensuremath{\mathbb{R}}}
\newcommand{\N}{\ensuremath{\mathbb{N}}}
\newcommand{\F}{\ensuremath{\mathbb{F}}}
\newcommand{\Pot}{\ensuremath{\mathcal{P}}}

\DeclareMathOperator{\reachOp}{E}
\newcommand{\E}[1]{\reachOp(#1)}


\newcommand{\rowvec}[2]{\begin{pmatrix}
#1&#2\\
\end{pmatrix}}

\newcommand{\colvec}[2]{\begin{pmatrix}
#1\\
#2
\end{pmatrix}}

\newcommand{\deactivateGlossaries}
{
    \renewcommand{\makenoidxglossaries}{}
    \renewcommand{\printnoidxglossaries}{}
}

%\deactivateGlossaries

\makenoidxglossaries

\usepackage{listings}
\lstset{basicstyle=\footnotesize, captionpos=b, breaklines=true, showstringspaces=false, tabsize=2, frame=lines, numbers=left, numberstyle=\tiny, xleftmargin=2em, framexleftmargin=2em}
% \makeatletter
% \def\l@lstlisting#1#2{\@dottedtocline{1}{0em}{1em}{\hspace{1,5em} Lst. #1}{#2}}
% \makeatother

\definecolor{javared}{rgb}{0.6,0,0} % for strings
\definecolor{javagreen}{rgb}{0.25,0.5,0.35} % comments
\definecolor{javapurple}{rgb}{0.5,0,0.35} % keywords
\definecolor{javadocblue}{rgb}{0.25,0.35,0.75} % javadoc
\definecolor{gray}{rgb}{0.6,0.6,0.6}
 
\lstset{
	language=Java,
	basicstyle=\ttfamily\footnotesize,
	keywordstyle=\color{javapurple}\bfseries,
	stringstyle=\color{javared},
	commentstyle=\color{javagreen}\itshape\bfseries,
	morecomment=[s][\color{javadocblue}]{/**}{*/},
	numbers=left,
	numberstyle=\tiny\color{gray},
	stepnumber=1,
	numbersep=10pt,
	tabsize=3,
	showspaces=false,
	showstringspaces=false
}

\lstdefinestyle{inline}{
	basicstyle=\ttfamily\normalsize,
	keywordstyle=\ttfamily\normalsize,
	breaklines=true,
	breakatwhitespace=true%,
	%literate={A}{}{0\discretionary{}{A}{A}}
  %         {B}{}{0\discretionary{}{B}{B}}
  %         {C}{}{0\discretionary{}{C}{C}}
  %         {D}{}{0\discretionary{}{D}{D}}
  %         {E}{}{0\discretionary{}{E}{E}}
  %         {F}{}{0\discretionary{}{F}{F}}
  %         {G}{}{0\discretionary{}{G}{G}}
  %         {H}{}{0\discretionary{}{H}{H}}
  %         {I}{}{0\discretionary{}{I}{I}}
  %         {J}{}{0\discretionary{}{J}{J}}
  %         {K}{}{0\discretionary{}{K}{K}}
  %         {L}{}{0\discretionary{}{L}{L}}
  %         {M}{}{0\discretionary{}{M}{M}}
  %         {N}{}{0\discretionary{}{N}{N}}
  %         {O}{}{0\discretionary{}{O}{O}}
  %         {P}{}{0\discretionary{}{P}{P}}
  %         {Q}{}{0\discretionary{}{Q}{Q}}
  %         {R}{}{0\discretionary{}{R}{R}}
  %         {S}{}{0\discretionary{}{S}{S}}
  %         {T}{}{0\discretionary{}{T}{T}}
  %         {U}{}{0\discretionary{}{U}{U}}
  %         {V}{}{0\discretionary{}{V}{V}}
  %         {W}{}{0\discretionary{}{W}{W}}
  %         {X}{}{0\discretionary{}{X}{X}}
  %         {Y}{}{0\discretionary{}{Y}{Y}}
  %         {Z}{}{0\discretionary{}{Z}{Z}}
}

\newcommand{\code}[1]{\lstinline[style=inline]!#1!}
%TODO fix overfull hboxes LATER
\newcommand{\class}[1]{\code{#1}}
\newcommand{\const}[1]{\code{#1}}
\newcommand{\var}[1]{\code{#1}}

\lstset{escapeinside={(*}{*)}}

\newcommand{\LSset}[2]{\scriptsize $\begin{aligned}&\{#1\}_L\\&\{#2\}_S\end{aligned}$}


\tikzset{
    transV/.style={transition, fill=black, minimum height = 12mm, minimum width = 1.5mm,inner sep = 0mm},
    transH/.style={transition, fill=black, minimum width = 12mm, minimum height = 1.5mm,inner sep = 0mm},
    node distance=1.5
}

\DeclareSIUnit{\bit}{bit}
\DeclareSIUnit{\byte}{B}


\newcommand{\todo}[1]{\fxnote{{\color{red}#1}}}
\newcommand{\TODO}[2]{\fxnote*{{\color{red}#1}}{\underline{\emph{#2}}}}

\newlength{\wrapfigwidth}

\usepackage[headsepline,footsepline]{scrlayer-scrpage}
\renewcommand{\sectionmark}[1]{\markboth{\Ifnumbered{section}{Kapitel \thesection}{}}{#1}}
\renewcommand{\subsectionmark}[1]{}
\renewcommand{\subsubsectionmark}[1]{}
\lohead{\leftmark}
\rohead{\rightmark}
\chead[]{}
\lofoot{Entwicklung einer nebenläufigen Architektur für eine 3D-Spielebibliothek}
\cfoot[]{}
\rofoot{Seite \pagemark}
\pagestyle{scrheadings}
\setkomafont{pageheadfoot}{\small}
\newglossaryentry{na}{
first={\emph{nicht-atomar}},
name={nicht-atomar},
plural={nicht-atomare},
user1={nicht-atomaren},
description={Beschreibung einer Anweisung, die selbst aus weiteren Anweisungen besteht}
}

\newacronym[]{fps}{FPS}{Frames per second (dt. Bilder pro Sekunde)}
\DeclareSIUnit{\fps}{\ac{fps}}
\newacronym[]{sot}{SoT}{System on a Thread}
\begin{document}
\thispagestyle{empty}
\includegraphics[scale=0.2]{IM_Logo.pdf}
\begin{center}
	\vspace*{1cm}
	\Large
	\textbf{Konzeption, Integration und Performanceanalyse einer nebenläufigen Architektur für eine 3D-Spielebibliothek}\\
	\vspace*{2cm}
	\large
	An der Fakultät für Informatik und Mathematik der\\
	Ostbayerischen Technischen Hochschule Regensburg\\
	im Studiengang\\
	Informatik\\[2em]
	eingereichte\\
	\vspace*{1cm}
	\Large
	\textbf{Masterarbeit}\\[2em]
	\large
	zur Erlangung des akademischen Grades des\\
	Master of Science (M.Sc.)
	\vspace*{1cm}
	\Large

	\vfill
	\normalsize
	\begin{tabular}{rl}
	    \rule{0mm}{1ex}\textbf{Vorgelegt von:} & Florian Loher \\
		\rule{0mm}{1ex}\textbf{Matrikelnummer:} & \hspace*{-0.5em}\begin{tabular}[t]{r}3252714\end{tabular} \\ 
		~\\
		\rule{0mm}{1ex}\textbf{Erstgutachter:} & Prof.\ Dr.\ Carsten Kern \\ 
		\rule{0mm}{1ex}\textbf{Zweitgutachter:} & Prof.\ Dr.\ Daniel Jobst \\[2em]
		\rule{0mm}{1ex}\textbf{Abgabedatum:} & 21.\ April 2022 \\ 
	\end{tabular} 
\end{center}
\clearpage
\null\thispagestyle{empty}\clearpage
\tableofcontents
\chapter{Einleitung}
\section{Blocklib}
\section{Nebenläufige Programmierung}
\section{Testen und Testframeworks}

\chapter{Grundlagen}
\subsection{Petri-Netze}\label{sec:petri}
Es gibt einige Ansätze zur Modellierung verteilter oder nebenläufiger Programme. Darunter finden sich auch Petri-Netze~\cite{Murata1989}. Mit einem Petrinetz kann ein modelliertes Programm mathematisch formal definiert und analysiert werden. Der Begriff Petri-Netz beschreibt auch eine Familie von verwandten Modellen. Das hier im Folgenden beschriebene Petri-Netz Modell wird auch Platz-Transitions-Netz (PT-Netz) genannt. Formal lässt sich ein PT-Netz $N$ als 5-Tupel $ N=(P,T,F,W,m_0)$ beschreiben. Dabei sind 
\begin{align*}
	&P  \quad \text{eine endliche Menge von Plätzen,}\\
	&T  \quad \text{eine endliche Menge von Transitionen,}\\
	&F \subseteq (P\times T) \cup (T \times P) \quad \text{die Menge der Relationen zwischen Plätzen und Transitionen,}\\
	&W: F \mapsto \N  \quad \text{eine Gewichtungsfunktion der Relationen und}\\
	&m_0: P \mapsto \N_0   \quad \text{die anfängliche Markierung der Plätze~\cite{Murata1989}.}\\
\end{align*}
Der durch ein PT-Netz beschriebene Graph kann auch grafisch dargestellt werden. Die Plätze $p \in P$ werden dabei als Kreise dargestellt, die Transitionen $ t \in T$ durch schwarz gefüllte Rechtecke, die Relationen $ f \in F$ durch gerichtete Kanten zwischen den Kreisen und Rechtecken, wobei die Elemente der Gewichtungsfunktion $W$ an die jeweiligen Kanten gesetzt werden. Die Markierung wird durch schwarze Punkte in den Plätzen dargestellt. Abbildun~\ref{fig:petrinet} zeigt ein simples Beispiel eines PT-Netzes  formal beschrieben (Abbildung~\ref{fig:petrinet:formal}) und seine grafische Repräsentation (Abbildung~\ref{fig:petrinet:graph}).
\begin{figure}
	\centering
	\begin{subfigure}[b]{.48\textwidth}
		$\begin{aligned}
			N &= (P,T,F,W,m_0)\\
			P &= \{p_1, p_2\}\\
			T &= \{t_1, t_2\}\\
			F &= \{(p_1, t_1), (p_2, t_2), (t_1, p_2)\}\\
			W &= \{((p_1, t_1),1), ((p_2, t_2),2), ((t_1, p_2), 1)\}\\
			m_0 &= \{(p_1, 1), (p_2, 0)\}
		\end{aligned}$
		\subcaption{}\label{fig:petrinet:formal}
	\end{subfigure}
	\hfill
	\begin{subfigure}[b]{.48\textwidth}
		\begin{tikzpicture}
			\node[place, label=$p_1$, tokens=1] (p1) at (0,0) {};
			\node[transV,label=$t_1$, right = of p1] (t1){};
			\node[place, label=$p_2$, right = of t1] (p2) {};
			\node[transV,label=$t_2$, right = of p2] (t2){};

			\draw 
			(p1) edge[post] node[above] {1} (t1)
			(t1) edge[post] node[above] {2} node[below=1cm] {$N$} (p2)
			(p2) edge[post] node[above] {1} (t2);
		\end{tikzpicture}
		\subcaption{}\label{fig:petrinet:graph}
	\end{subfigure}
	\caption{}\label{fig:petrinet}
\end{figure}

\subsubsection{Transitionsregel} Um nun das Verhalten von Programmen beschreiben zu können, kann die Markierung eines Petrinetzes anhand der folgenden sogenannten Transitions-Regel geändert werden~\cite{Murata1989}.
\begin{enumerate}
	\item Eine Transition $t_i \in T$ heißt \emph{aktiviert}, wenn für alle $(p,t_i) \in F $ gilt $ W((p,t_i)) \leq M(p)$ wobei $M$ die aktuelle Markierung ist. Es müssen also in allen Plätzen mit eingehenden Kanten genügend Marken bezüglich der Gewichtungsfunktion existieren. Diese Plätze werden auch \emph{Vorbedinungen} für die Transition genannt. Wenn ein Platz genügend Marken für eine Transition enthält, nennt man die Vorbedinung (bezüglich dieser Transition) \emph{erfüllt}.
	\item Eine aktivierte Transition kann, muss aber nicht \emph{feuern}.
	\item Feuert eine aktivierte Transition $t_i$, geht eine Markierung $M$ in eine Markierung $M'$ über. Dies erfolgt nach den folgenden Regeln.
	\begin{enumerate}
		\item $P$ wird in die disjunkten Mengen $P_\emptyset, P_\rightarrow, P_\leftarrow, P_\leftrightarrow$ unterteilt, wobei 
		\item $P_\emptyset = \{p \in P | (p,t_i) \notin F \land (t_i,p) \notin F\}$ die unbeteiligten Plätze sind,
		\item $P_\rightarrow = \{p \in P | (p,t_i) \in F \land (t_i,p) \notin F\}$ die Plätze mit eingehenden Kanten zur Transition,
		\item $P_\leftarrow = \{p \in P | (p,t_i) \notin F \land (t_i,p) \in F\}$ die Plätze mit ausgehenden Kanten aus der Transition und 
		\item $P_\leftrightarrow = \{p \in P | (p,t_i) \in F \land (t_i,p) \in F\}$ die Plätze mit eingehenden und ausgehenden Kanten.
		\item Dann gilt $
			M'(p) = \left\{ 
				\begin{aligned}
					& M(p) && \; , \; p \in P_\emptyset\\
					& M(p)-W((p,t_i)) && \; , \; p \in P_\rightarrow\\
					&M(p)+W((t_i,p)) && \; , \; p \in P_\leftarrow\\
					& M(p)-W((p,t_i))+W((t_i,p)) && \; , \; p \in P_\leftrightarrow
				\end{aligned}
				\right\}
		$
	\end{enumerate}
\end{enumerate}
Betrachtet man Abbildung~\ref{fig:petrinet} erneut, lässt sich erkennen, dass Transition $t_1$ aktiviert ist und Transition $t_2$ nicht, da in $p_2$ keine Marken sind.

\subsubsection{Erreichbarkeitsgraph}
Mittels der Transitionsregel lassen sich in einem Petri-Netz $N$ ausgehend von der Anfangsmarkierung $m_0$ gegebenenfalls weitere Markierungen erzeugen. Durch die wiederholte Anwendung der Transitionsregel auf alle daraus entstehenden Markierungen lässt sich ein Graph erstellen, der alle von der Anfangsmarkierung aus erreichbaren Markierungen erhält, den \emph{Erreichbarkeitsgraph} $\E{N}$ des Petri-Netzes.
\begin{figure}
	\begin{tikzpicture}[node distance=5mm,auto,state/.append style={rounded rectangle}]
		\node[state, initial] (q1) {$\{(p_1,1),(p_2,0)\}$};
		\node[state] (q2) [right=of q1] {$\{(p_1,0),(p_2,2)\}$};
		\node[state] (q3) [right=of q2] {$\{(p_1,0),(p_2,1)\}$};
		\node[state] (q4) [right=of q3] {$\{(p_1,0),(p_2,0)\}$};
		\path[->] 
		(q1) edge node {$t_1$} (q2)
		(q2) edge node {$t_2$} node[below=6mm] {Graph mit vollständigen Knotenbezeichnungen}(q3)
		(q3) edge node {$t_2$} (q4);

		\node[state,initial] (q5) [below=1.2cm of q1]{$(1,0)$};
		\node[state] (q6) [below=1.2cm of q2] {$(0,2)$};
		\node[state] (q7) [below=1.2cm of q3] {$(0,1)$};
		\node[state] (q8) [below=1.2cm of q4] {$(0,0)$};
		\path[->] 
		(q5) edge node {$t_1$} (q6)
		(q6) edge node {$t_2$} node[below=6mm] {Graph mit gekürzten Knotenbezeichnungen}(q7)
		(q7) edge node {$t_2$} (q8);
		\end{tikzpicture}
		\caption{Erreichbarkeitsgraph des Petri-Netzes aus Abbildung~\ref{fig:petrinet}. Der Graph ist zweimal dargestellt. Oben enthalten die Knoten des Graphen die vollständige Auszeichnung der Markierung, unten wird die Kurzschreibweise für die Knotenbezwichnungen genutzt, in der die Position den Platz kodiert. Der \enquote{start} Pfeil kennzeichnet die anfängliche Markierung.}\label{fig:reachability}
\end{figure}
In Abbildung~\ref{fig:reachability} wird der Erreichbarkeitsgraph des in Abbildung~\ref{fig:petrinet} dargestellten Petri-Netzes gezeigt. Die Knoten des Graphen sind die Markierungen des Petri-Netzes, die durch die Transitionsregel erreicht werden können. Da die Plätze des Petri-Netzes nummeriert sind, kann die Markierung verkürzt geschrieben werden, indem der Index des Tupels der Knotenbezeichnung den Platz beschreibt und die Zahl die Anzahl der Marken des Platzes kennzeichnet. In der Abbildung ist oben der Graph mit den vollständigen Markierungsbezeichnungen dargestellt. Unten wird die verkürzte Schreibweise genutzt.

ist die Anzahl der Marken pro Platz in jeder Markierung des Erreichbarkeitsgraphen maximal $1$, kann die Schreibweise weiter verkürzt werden, indem nur die Plätze beziehungsweise die Indizes der Plätze, die eine Marke enthalten genannt werden. Abbildung~\ref{fig:1-bpetrinet} zeigt eine Abwandlung des Petri-Netzes aus Abbildung~\ref{fig:petrinet}, das der gerade beschriebenen Eigenschaft entspricht, und den zugehörigen Erreichbarkeitsgraphen in verkürzter Schreibweise.
\begin{figure}
	\centering
	\begin{subfigure}[b]{.48\textwidth}
		\begin{tikzpicture}
			\node[place, label=$p_1$, tokens=1] (p1) at (0,0) {};
			\node[transV,label=$t_1$, right = of p1] (t1){};
			\node[place, label=$p_2$, right = of t1] (p2) {};
			\node[transV,label=$t_2$, right = of p2] (t2){};

			\draw 
			(p1) edge[post] node[above] {1} (t1)
			(t1) edge[post] node[above] {1} (p2)
			(p2) edge[post] node[above] {1} (t2);
		\end{tikzpicture}
		\subcaption{}\label{fig:1-bpetrinet:graph}
	\end{subfigure}
	\hfill
	\begin{subfigure}[b]{.48\textwidth}
		\begin{tikzpicture}[node distance=5mm,auto,state/.append style={rounded rectangle}]
			\node[state, initial] (q1) {$1$};
			\node[state] (q2) [right=of q1] {$2$};
			\node[state] (q3) [right=of q2] {};
			\path[->] 
			(q1) edge node {$t_1$} (q2)
			(q2) edge node {$t_2$} (q3);
		\end{tikzpicture}
		\subcaption{}\label{fig:1-bpetrinet:reachability}
	\end{subfigure}
\caption{}\label{fig:1-bpetrinet}
\end{figure} 

\subsubsection{Weitere Definitionen}
Um die Arbeit mit Petri-Netzen zu vereinfachen, werden nun noch einige weitere Definitionen eingeführt.

\begin{enumerate}
	\item Die Menge der Knoten mit Kanten zu einem Knoten $k \in P\cup T$ heißt \emph{Vorbereich} des Knotens $k$ und ist definiert als $^\circ k = \{v | (v,k) \in F\}$.
	\item Die Menge der Knoten mit Kanten von einem Knoten $k \in P\cup T$ heißt \emph{Nachbereich} des Knotens $k$ und ist definiert als $k^\circ  = \{n | (k,n) \in F\}$.
	\item Die Menge der Knoten einer Markierung $m$ mit mindesten $n$ Marken $\{x | m(x)\geq n\}$ wird $M_{\geq n}$ genannt.
	\item Das Maximum einer Markierung $m$, $\deg(m) \coloneqq \max \{m(p)|p\in P\}$, wird ihr \emph{Grad} genannt.
	\item Ein Petri-Netz $N$ heißt \emph{$n$-beschränkt}, wenn $n$ das Maximum der Grade der Markierungen des Erreichbarkeitsgraphen $(V,K)=\E{N}$ ist, also $n = \max\{\deg(m)| m\in V \}$.
	\item Gilt für eine Kante $((a, b),w) \in W$, dass $w = 1$, so kann im dazugehörigen Graphen die Beschriftung entfallen. Kanten ohne Beschriftung haben also ein implizites Gewicht von $1$.
\end{enumerate}



\subsubsection{Erweiterte Petri-Netze}
Um das Konzept von Variablenzugriffen einfacher zu modellieren führen \textcite{Goel1990} ein erweitertes Petri-Netz Modell ein. Dabei wird das 5-Tupel des Petri-Netzes um die folgenden Komponenten erweitert:
\begin{enumerate}
	\item eine Menge von Variablen $V$,
	\item eine Funktion $L: T \mapsto \Pot(V)$, die \emph{lesenden Zugriff} auf die Variablen modelliert, und
	\item eine Funktion $S: T \mapsto \Pot(V)$, die \emph{schreibenden Zugriff} auf die Variablen modelliert.
\end{enumerate}
$\Pot(V)$ ist dabei die Potenzmenge von $V$, also die Menge aller Teilmengen von $V$.
Ein Beispiel für ein erweitertes Petri-Netz, das ansonsten identisch zu dem Petri-Netz in Abbildung~\ref{fig:petrinet} ist, ist in Abbildung~\ref{fig:augpetrinet} gegeben. Die Formale Definition zeigt die neuen Mengen 
\begin{figure}
\centering
	\begin{minipage}[c]{.49\textwidth}
		$\begin{aligned}
			N &= (P,T,F,W,m_0, V, L, S)\\
			P &= \{p_1, p_2\}\\
			T &= \{t_1, t_2\}\\
			F &= \{(p_1, t_1), (p_2, t_2), (t_1, p_2)\}\\
			W &= \{((p_1, t_1),1), ((p_2, t_2),2), ((t_1, p_2), 1)\}\\
			m_0 &= \{(p_1, 1), (p_2, 0)\}\\
			V &= \{a,b\}\\
			L &= \{(t_1,\{a,b\}), (t_2,\emptyset)\}\\
			S &= \{(t_1,\{a\}), (t_2,\emptyset)\}
		\end{aligned}$
	\end{minipage}
	\hfill
	\begin{minipage}[c]{.49\textwidth}
		\begin{tikzpicture}
			\node[place, label=$p_1$, tokens=1] (p1) at (0,0) {};
			\node[transV,label=$t_1$, right = of p1, label=below:{\LSset{a,b}{a}}] (t1){};
			\node[place, label=$p_2$, right = of t1] (p2) {};
			\node[transV,label=$t_2$, right = of p2] (t2){};

			\draw 
			(p1) edge[post] node[above] {1} (t1)
			(t1) edge[post] node[above] {2} node[below=2cm] {$N$} (p2)
			(p2) edge[post] node[above] {1} (t2);
		\end{tikzpicture}
	\end{minipage}
	\caption{Beispiel eines erweiterten Petri-Netzes. Links ist die formale Definition gegeben, rechts der dazugehörige Graph. In Platz $p_2$ wird auf die Variablen $a$ und $b$ lesend zugegriffen und auf Variable $a$ schreibend.}\label{fig:augpetrinet}
\end{figure}
Nachdem nun die formalen Grundlagen gelegt sind, werden in diesem Abschnitt Grundbegriffe eingeführt, die notwendig für das Verständnis der nächsten Kapitel sind. Besonders wichtig sind insbesondere die Begriffe \enquote{Thread}, \enquote{Nebenläufigkeit} und \enquote{Wettkampfbedingung}, da diese Begriffe in der Arbeit intensiv behandelt werden.


\subsubsection{Definition von Thread}
Um definieren zu können, was ein Thread ist, müssen zuerst einige weitere Grundbegriffe eingeführt werden. 

Als \gls{Programm} wird in dieser Arbeit die konkrete Niederschrift eines Algorithmus bezeichnet. Die Elemente, aus denen das \gls{Programm} besteht, werden als \glspl{Anweisung} bezeichnet. Bei der Ausführung eines \glsuseri{Programm} werden Einzelschritte durchlaufen, diese bezeichnet man als \glspl{Aktivitaet}. Eine Sequenz von \glspl{Aktivitaet}, die ein isoliertes Problem abarbeitet, wird \emph{Prozess} genannt. Um diese Definition des Prozessbegriffs von späteren Definitionen abzugrenzen wird er im Folgenden \gls{Rechenprozess} genannt. Die Ausführungseinheit, auf der die Schritte eines \glsuseri{Rechenprozess} durchgeführt werden, wird \emph{Prozessor} genannt~\cite[S.~22]{Herrtwich1989}. In der Regel besitzt ein Rechner, beispielsweise ein PC, mehrere Ausführungseinheiten. Bei PCs gibt es eine Recheneinheit, die üblicherweise die meisten \glspl{Rechenprozess} ausführt und andere Prozessoren koordiniert, den \ac{cpu}.

Wenn ein \gls{Programm} auf einem Betriebssystem ausgeführt wird, erzeugt das Betriebssystem einen isolierten Adressraum, in dem das \gls{Programm} ausgeführt wird. Ein \gls{Rechenprozess}, der auf diese Weise ausgeführt wird, wird in dieser Arbeit als \emph{(System-)Prozess} bezeichnet. \textcite[Kapitel~2]{Tanenbaum2016} liefern eine gute Übersicht über Systemprozesse, dort als Prozesse bezeichnet. Systemprozesse können selbst weitere Systemprozesse starten, wenn das Betriebssystem dies zulässt. Diese Systemprozesse besitzen dann ihren eigenen Adressraum. Die Anzahl der Systemprozesse, die auf einem Rechner laufen, ist meist höher als die Anzahl der Prozessoren des Rechners. Somit können nicht alle Prozesse zur gleichen Zeit ausgeführt werden. Damit dennoch alle Prozesse voranschreiten können, wechseln moderne Betriebssysteme die Systemprozesse, die auf den Prozessoren des Rechners ausgeführt werden, in schneller Folge. Dabei muss das Betriebssystem einige zu den Systemprozessen gehörende Daten, den \emph{Prozesskontrollblock}, tauschen, sodass der jeweils gerade ausführende Systemprozess seinen eigenen Prozesskontrollblock zur Verfügung hat. Dieser Tausch wird \emph{Kontextwechsel} genannt~\cite[S.~59]{Tanenbaum2016}.

Das Speichern und Laden der Prozesskontrollblöcken im Rahmen der Kontextwechsel von Systemprozessen benötigt eine nicht zu vernachlässigende Menge an Zeit. Dieser Zeitverbrauch wird auch als \emph{Overhead} bezeichnet. Um den Overhead von Kontextwechseln zu vermindern, bieten moderne Betriebssysteme Systemprozessen die Möglichkeit \glspl{Rechenprozess} zu erzeugen, die denselben Adressraum wie der erzeugende Systemprozess besitzen. Ein auf diese Art erzeugter \gls{Rechenprozess} heißt \emph{Thread}. Da die Menge der thread-eigenen Daten, des \emph{Threadkontrollblocks}, deutlich geringer ist als die des Prozesskontrollblocks, erzeugt der Kontextwechsel zwischen Threads einen geringeren Overhead. Zudem können Threads aufgrund des gemeinsamen Adressraums einfacher auf geteilte Ressourcen zugreifen. Das vereinfacht die Kooperation zwischen diesen \glsuserii{Rechenprozess}~\cite[S.~139~ff.]{Tanenbaum2016}.

\paragraph{Anweisungen und Petri-Netze}
Um formal Eigenschaften eines \glsuseri{Programm} zu beschreiben, ist es möglich dieses mittels eines Petri-Netzes (automatisiert) zu modellieren. Dazu müssen Petri-Netz-Konstrukte genutzt werden, die die \glspl{Anweisung} des \glsuseri{Programm} abbilden können. Für die Analyse nebenläufiger \glspl{Programm} sind nur sogenannte \emph{Rendezvous}-\glspl{Anweisung} sowie Kontrollstrukturen relevant~\cite{Goel1990}. Für diese lassen sich Unter-Petri-Netze definieren, die deren Verhalten modellieren. Eine Anweisung entspricht dann einem bestimmten Unter-Petri-Netze. Eine Auflistung für die Modellierung nebenläufiger \glspl{Programm} benötigter Unter-Petri-Netze ist in \cite{Goel1990} zu finden.  Da die Modellierung von Verzweigungen durch Unter-Petri-Netze in \cite[Abbildung 3.1]{Goel1990} semantisch nicht ganz korrekt ist, wird sie in Abbildung~\ref{fig:supnetifelse} korrigiert dargestellt. 
\begin{figure}
	\centering
	\begin{tikzpicture}[node distance=2cm,on grid, auto]
		\node[place, label=$p_1$] (p1) {};
		\node[transV, label=\texttt{if}, above right = of p1] (if) {};
		\node[transV, label=\texttt{else},below right = of p1] (else) {};
		\node[right = of if] (dots1){$\dots$};
		\node[right = of else] (dots2){$\dots$};
		\node[place, label=$p_2$, right = of dots1] (p2) {};
		\node[place, label=$p_3$, right = of dots2] (p3) {};
		\node[transV, label=$t_1$, right = of p2] (t1) {};
		\node[transV, label=$t_2$, right = of p3] (t2) {};
		\node[place, label=$p_4$, below right = of t1] (p4) {};
		\node[transV, label=\texttt{end if}, right = of p4] (endif) {};
	
	
	
		\draw 
		(p1) edge[post] (if)
		(p1) edge[post] (else)
		(if) edge[post] (dots1)
		(else) edge[post] (dots2)
		(dots1) edge[post] (p2) 
		(dots2) edge[post] (p3) 
		(p2) edge[post] (t1)
		(p3) edge[post] (t2)
		(t1) edge[post] (p4)
		(t2) edge[post] (p4)
		(p4) edge[post] (endif)
		;
	\end{tikzpicture}
	\caption[Semantisch korrigiertes Unter-Petri-Netz zu Modellierung von Verzweigungen.]{Semantisch korrigiertes Unter-Petri-Netz zu Modellierung von Verzweigungen nach \cite[Abbildung~3.1]{Goel1990}.}\label{fig:supnetifelse}
\end{figure}

Die Plätze $p_1,p_2,p_3,p_4$ und die Transitionen $t_1,t_2$ sind Hilfselemente die die semantische Korrektheit der Transitionen \code{if}, \code{else} und \code{end if} sicherstellen. Eine Markierung kann exklusiv nur von \code{if} oder von \code{else} zum Feuern verwendet werden. Durch die Transitionen $t_1$ und $t_2$ entsteht eine Markierung in $p_5$ wodurch \code{end if} feuern kann egal ob anfangs \code{if} oder \code{else} gefeuert hat. $p_2$ und $p_3$ existierten, damit das Petri-Netz ein bipartiter Graph bleibt, sich also immer Plätze und Transitionen \enquote{abwechseln}. Die Modellierung eines \code{switch}-Statements erfolgt analog, indem die Anzahl der Pfade zwischen $p_1$ und $p_4$ angepasst wird.

\subsubsection{Nebenläufigkeit}\label{sec:nebenl}
Im vorherigen Abschnitt wurde beschrieben, dass verschiedene \glspl{Rechenprozess} unabhängig von einander \enquote{gleichzeitig} ablaufen können. Dieses \enquote{gleichzeitig} wird jetzt genauer beleuchtet. Ein \gls{Programm} wird häufig als \emph{Sequenz} von \glspl{Anweisung} verstanden. In der Definition von \gls{Programm} in dieser Arbeit wird bewusst auf das Wort Sequenz verzichtet, denn bei vielen \glsuserii{Programm} kann es bei bestimmten \glspl{Anweisung} irrelevant sein, in welcher Reihenfolge oder ob sie sogar gleichzeitig ausgeführt werden. Man betrachte das folgende \gls{Programm} in Listing~\ref{lst:squareSeqEx}, das die Quadratsumme zweier Ganzzahlen (im Folgenden auch Quadratsumme genannt) ausgibt.

\begin{lstlisting}[caption={[Beispiel eines \glsentryuseri{Programm} das sequentiell die Summe von Quadraten zweier Ganzzahlen berechnet.]Beispiel eines \glsuseri{Programm} das die Summe von Quadraten zweier Ganzzahlen berechnet. Die Berechnung der Quadratzahlen wird nacheinander in einer fest definierten Sequenz durchgeführt.}, label={lst:squareSeqEx},float={!htbp}]
printSummedSquare(int a, int b){
  x <- a*a (*\label{lst:squareSeqEx:aa}*)
  y <- b*b (*\label{lst:squareSeqEx:bb}*)
  print(x+y) (*\label{lst:squareSeqEx:print}*)
}
\end{lstlisting}
Es ist leicht zu erkennen, dass es vollkommen egal ist, ob zuerst \code{x} oder \code{y} ausgerechnet wird oder die Berechnungen simultan stattfinden. Einzig der Ausgabebefehl in Zeile~\ref{lst:squareSeqEx:print} darf erst ausgeführt werden, nachdem \code{x} und \code{y} berechnet wurden. Zuzulassen, dass die Ausführreihenfolge in dem \gls{Programm} nicht definiert ist, führt dazu, dass das \gls{Programm} präziser der Problembeschreibung \enquote{gib die Quadratsumme zweier Ganzzahlen aus} entspricht, und ermöglicht auch eine potenziell schnellere Berechnung, da die Quadrate gleichzeitig berechnet werden können.

Paare oder Gruppen von \glspl{Anweisung}, die gleichzeitig oder in beliebiger Reihenfolge ausgeführt werden, heißen \emph{nebenläufig}. Da die Ausführungsreihenfolge der Anweisungen in Zeile~\ref{lst:squareSeqEx:aa} und Zeile~\ref{lst:squareSeqEx:bb} egal ist, dürfen diese Anweisungen auch nebenläufig definiert sein.

Eine der obigen Problemstellung, \enquote{gib die Quadratsumme zweier Ganzzahlen aus}, nähere Implementierung der Quadratsumme könnte unter Nutzung von Nebenläufigkeit wie Listing~\ref{lst:squareConcEx} aussehen. Dabei wird auf die Schreibweise von \textcite[S.~16]{Herrtwich1989} zurückgegriffen.
\begin{lstlisting}[caption={[Beispiel eines \glsuseri{Programm} mit nebenläufigem Code in einem \code{conc}-Block.]Beispiel eines \glsuseri{Programm} mit nebenläufigem Code in einem \code{conc}-Block. Das \gls{Programm} gibt die Summe von zwei Quadratzahlen aus, wobei die Berechnung der Quadratzahlen nebenläufig stattfindet.}, label={lst:squareConcEx},float={!htbp}]
printSummedSquare(int a, int b){
  conc 
    x <- a*a ||
    y <- b*b
  conc end
  print(x+y)
}
\end{lstlisting}
In einem \code{conc}-Block, definiert durch \code{conc} und \code{conc end}, sind alle \glspl{Anweisung}, die durch \code{||} getrennt sind, als nebenläufig zu verstehen. 

Wie man an dem Beispiel in Listing~\ref{lst:squareConcEx} erkennen kann, beschäftigt sich Nebenläufigkeit besonders mit der Struktur und der Möglichkeit der Zusammensetzung voneinander unabhängiger \glspl{Anweisung}. Somit ist es Aufgabe der nebenläufigen Programmierung, Probleme oder Aufgaben in unabhängige Teile zu zerlegen und zu strukturieren~\cite{Pike2012,Hettel2016}. Blöcke von \glspl{Anweisung}, die (bezüglich der Aufgabe) nebenläufig sein könnten, aber als Anweisungssequenz definiert sind (siehe Listing~\ref{lst:squareSeqEx} Zeile 2 und 3), werden in dieser Arbeit als \emph{pseudo-sequentialisiert} bezeichnet.

Bei der Definition nebenläufiger Anweisungen ist allerdings Vorsicht geboten, da sichergestellt werden muss, dass zwei Anweisungen auch tatsächlich nebenläufig sein dürfen. Eine \gls{Anweisung} kann beispielsweise durch einen Compiler in mehrere \glspl{Anweisung} geteilt werden oder selbst schon aus mehreren \glspl{Anweisung} bestehen (man denke zum Beispiel an Funktionsaufrufe). Solche \glspl{Anweisung} werden \gls{nichtAtomar} genannt. Eine Menge von nebenläufigen \glspl{Anweisung}, von denen mindestens eine \gls{nichtAtomar} ist, kann \emph{verzahnt} ausgeführt werden. Eine Ausführung von \glspl{Anweisung} ist verzahnt, wenn zwischen der Ausführung der Teile einer \glsuseri{nichtAtomar} \gls{Anweisung} andere \glspl{Anweisung} ausgeführt werden. Abbildung \ref{fig:concAnweisungen} zeigt zur Veranschaulichung die verschiedenen Möglichkeiten, wie zwei nebenläufige \glspl{nichtAtomar} \glspl{Anweisung} im Zeitverlauf ausgeführt werden können, sowohl auf einem Prozessor als auch auf zwei Prozessoren. Dadurch kann es passieren, dass Anweisungen, die zuerst den Anschein haben, nebenläufig sein zu können, nicht nebenläufig sein dürfen, weil eine verzahnte Ausführung der zugehörigen Aktivitäten ein unerwünschtes Ergebnis liefern würde. 

\begin{figure}[hbt]
	\centering
\newlength\aOne
\newlength\aTwo
\newlength\aThree
\newlength\bOne
\newlength\bTwo
\pgfmathsetlength{\aOne}{2cm}
\pgfmathsetlength{\aTwo}{2.25cm}
\pgfmathsetlength{\aThree}{2.5cm}
\pgfmathsetlength{\bOne}{2.75cm}
\pgfmathsetlength{\bTwo}{3cm}
\begin{subfigure}{\textwidth}
	\centering
\begin{tikzpicture}

	\draw[->] (0,0) -- (13,0) node[right] {Zeit};
	
	\node[draw,fill=red, anchor=south west, minimum width = \aOne] at (0, .5) (A1) {A Teil1};

	\node[draw,fill=red, anchor=west, minimum width = \aTwo] at (A1.east) (A2) {A Teil2};
	\node[draw,fill=red, anchor=west, minimum width = \aThree] at (A2.east) (A3) {A Teil3};
	
	\node[draw,fill=cyan, anchor=south west, minimum width = \bOne] at (0, 1.5) (B1) {B Teil1};
	\node[draw,fill=cyan, anchor=west, minimum width = \bTwo] at (B1.east) (B2) {B Teil2};

	\node[anchor=east] at (B1.west) {Prozessor 1};
	\node[anchor=east] at (A1.west) {Prozessor 2};
\end{tikzpicture}
\subcaption{Gleichzeitige Ausführung der \glspl{Anweisung} A und B auf zwei Prozessoren}
\end{subfigure}
\\[1.5em]
\begin{subfigure}{\textwidth}
	\centering
\begin{tikzpicture}

	\draw[->] (0,0) -- (13,0) node[right] {Zeit};
	
	\node[draw,fill=cyan, anchor=south west, minimum width = \bOne] at (0, .5) (B1) {B Teil1};
	\node[draw,fill=cyan, anchor=west, minimum width = \bTwo] at (B1.east) (B2) {B Teil2};
	\node[draw,fill=red, anchor=west, minimum width = \aOne] at (B2.east) (A1) {A Teil1};
	\node[draw,fill=red, anchor=west, minimum width = \aTwo] at (A1.east) (A2) {A Teil2};
	\node[draw,fill=red, anchor=west, minimum width = \aThree] at (A2.east) (A3) {A Teil3};

	\node[anchor=east] at (B1.west) {Prozessor 1};
\end{tikzpicture}
\subcaption{Sequenzielle Ausführung der \glspl{Anweisung} A und B auf einem Prozessor}
\end{subfigure}
\\[1.5em]
\begin{subfigure}{\textwidth}
	\centering
\begin{tikzpicture}

	\draw[->] (0,0) -- (13,0) node[right] {Zeit};
	
	\node[draw,fill=cyan, anchor=south west, minimum width = \bOne] at (0, .5) (B1) {B Teil1};
	\node[draw,fill=red, anchor=west, minimum width = \aOne] at (B1.east) (A1) {A Teil1};
	\node[draw,fill=red, anchor=west, minimum width = \aTwo] at (A1.east) (A2) {A Teil2};
	\node[draw,fill=cyan, anchor=west, minimum width = \bTwo] at (A2.east) (B2) {B Teil2};
	\node[draw,fill=red, anchor=west, minimum width = \aThree] at (B2.east) (A3) {A Teil3};

	\node[anchor=east] at (B1.west) {Prozessor 1};
\end{tikzpicture}
\subcaption{Verzahnte Ausführung der \glspl{Anweisung} A und B auf einem Prozessor}
\end{subfigure}

\caption[Mögliche Ausführungen nebenläufiger \glsentryplural{Anweisung}.]{Mögliche Ausführungen nebenläufiger \glspl{Anweisung}. Die Grafik bildet eine Abbildung aus \cite{Herrtwich1989} nach.}\label{fig:concAnweisungen}
\end{figure}

Werden nebenläufige \glspl{Anweisung} parallel ausgeführt, spricht man auch von \emph{Multiprocessing}. Handelt es sich bei den Prozessen um Threads, wird der Begriff \emph{Multithreading} genutzt.

\paragraph{Nebenläufigkeit und Petri-Netze}
Um ein eindeutiges Verständnis des hier genutzten Begriffs der Nebenläufigkeit zu geben, wird dieser auch in Bezug auf Petri-Netze definiert. Wichtig anzumerken ist, dass der hier beschriebene Begriff von dem üblichen Verständnis der Nebenläufigkeit in Petri-Netzen abweicht. 

Normalerweise werden Petri-Netze als inhärent nebenläufig verstanden, da durch die Transitionsregel zu jeder Zeit jede beliebige Transition feuern kann, deren Vorbedingungen erfüllt sind. Der hier verwendete Begriff bezieht sich allerdings darauf, dass \glspl{Anweisung} unabhängig voneinander, ohne sich gegenseitig zu beeinflussen, ausführbar sind. Das Feuern einer Transition kann aber den Vorbereich einer anderen Transition verändern, wodurch sie beeinflusst wird.

In dieser Arbeit heißen zwei Transitionen $t_1, t_2$ eines Petri-Netzes $N$ genau dann nebenläufig, wenn ${}^\circ t_1 \cap {}^\circ t_2 = \emptyset$ gilt und eine Markierung $m_\text{conc}$ im Erreichbarkeitsgraphen $\E{N}$ existiert, in der sowohl $t_1$ als auch $t_2$ aktiviert sind. Die Vorbereiche der Transitionen dürfen sich also nicht überschneiden.

\subsubsection{Folgen von Nebenläufigkeit}\label{sec:nebenl-folgen}
Wie in Abschnitt \ref{sec:nebenl} beschrieben, muss ein \gls{Programm} keine Sequenz von \glspl{Anweisung} sein, sondern kann auch nebenläufige \glspl{Anweisung} enthalten. Diese Möglichkeit hat eine Reihe von Folgen, die es bei der nebenläufigen Programmierung zu beachten gilt.
\paragraph{Nichtdeterminismus}
Enthält ein \gls{Programm} nebenläufige \glspl{Anweisung}, ist die Reihenfolge der daraus resultierenden \glspl{Aktivitaet} nicht definiert und kann sich bei jedem Programmdurchlauf ändern. Die Ausführung ist also \emph{nichtdeterministisch}, als direkte Folge der Nebenläufigkeit~\cite[S.~17~f.]{Herrtwich1989}. Erwartet man von einem \gls{Programm} \emph{Determiniertheit}\footnote{Es kann durchaus sein, dass Determiniertheit in einem \gls{Programm} nicht gewünscht ist. Man betrachte beispielsweise ein \gls{Programm}, das einen echten Zufallsgenerator beschreibt. Hier wäre Determiniertheit ein direkter Widerspruch zur Aufgabe des \glsuseri{Programm}.}, also die Eigenschaft, dass gleiche Eingaben immer zu den gleichen Ausgaben führen, ist Nichtdeterminismus in der Regel zu vermeiden. Determiniertheit kann aber auch bei einem Verzicht auf Determinismus sichergestellt werden. Liefert ein \gls{Programm} bei jeder beliebigen Ausführreihenfolge dasselbe Ergebnis, ist es trotz Nichtdeterminismus weiterhin determiniert, weil die Ausführreihenfolge für das Ergebnis keine Rolle spielt~\cite[S.~18~f.]{Herrtwich1989}. 
\paragraph{Nichtreproduzierbarkeit}
Da das Wissen und die Kontrolle über die Ausführreihenfolge abgegeben wird, wird die Nachvollziehbarkeit des Programmablaufs erschwert~\cite[S.~20]{Herrtwich1989}. Tritt beispielsweise ein Fehler auf, kann im Nachhinein nicht ermittelt werden, was die Ausführreihenfolge war, die zu dem Fehler geführt hat. Dasselbe Problem ergibt sich ebenfalls, wenn ein nichtdeterminiertes \gls{Programm} eine Lösung ausgibt und nachvollzogen werden soll, welche Ausführreihenfolge zu diesem Ergebnis geführt hat. Insbesondere das Testen und Debugging von Software wird dadurch erschwert, da es unmöglich ist bei jedem Test dieselben Bedingungen herzustellen, sodass beispielsweise Tests Fehler nicht verlässlich aufzeigen, weil diese nur bei bestimmten Ausführreihenfolgen auftreten~\cite[S.~20]{Herrtwich1989}. Somit muss ein Test entweder alle möglichen Ausführreihenfolgen simulieren oder es muss anderweitig sichergestellt werden, dass die Ausführreihenfolge der nebenläufigen \glspl{Anweisung} keine Rolle spielt.
\paragraph{Wettkampfbedingungen}
Ressourcen (zum Beispiel Drucker, Variablen und Dateien) und insbesondere Daten spielen in \glsuserii{Programm} eine zentrale Rolle. Greift eine Gruppe von nebenläufigen \glspl{Anweisung} auf dieselben (geteilten) Daten zu, ist die Reihenfolge der Zugriffe, wie alle anderen \glspl{Aktivitaet} der nebenläufigen \glspl{Anweisung}, beliebig. Wenn mindestens eine der \glspl{Anweisung} schreibend auf die geteilten Daten zugreift (diese also verändert), hängt der Zustand der Ausgabe des \glsuseri{Programm} im Allgemeinen von der Reihenfolge der Ausführung ab. Diese Situation wird \emph{Wettkampfbedingung} (engl. Race Condition) genannt~\cite{Hettel2016}. 

Erwähnenswert ist hier, dass Wettkampfbedingungen nur auftreten können, wenn die geteilten Daten nebenläufig \emph{geändert} werden. Ausschließlich lesende Zugriffe sind unkritisch, da die Daten unabhängig von der Ausführreihenfolge immer identisch sind. Wenn Daten von \glspl{Anweisung} zu jederzeit in einem Zustand aufgefunden werden, der korrekt ist, werden sie als \emph{threadsicher} bezeichnet.

Formal können Wettkampfbedingungen auch mittels erweiterter Petri-Netze definiert werden. Dazu betrachtet man, ob Variablen des erweiterten Petri-Netzes sich in den Lese- und Schreibmengen von nebenläufigen Transitionen überschneiden. Diese Variablen werden \emph{kritische} Variablen genannt. Die Menge der kritischen Variablen, die zur Existenz von Wettkampfbedingungen führen, ist gegeben durch
\begin{align*}
	V_\text{krit} = \;\bigcup_{\mathclap{\substack{s,t\, \in T\\s, t \text{ nebenläufig}}}} \;\left( \strut{S(s) \cap (S(t) \cup L(t))}\right).
\end{align*}

Ein Paar von Transitionen $(t_1, t_2)$ eines erweiterten Petri-Netzes \emph{erzeugt} eine Wettkampfbedingung, wenn die Transitionen dazu führen, dass eine Variable in die Menge der kritischen Variablen aufgenommen wird. Eine Markierung $m$ des Erreichbarkeitsgraphen $\E{N}$ eines Petri-Netzes $N$ enthält eine Wettkampfbedingung, wenn ein Paar von Transitionen $(t_1, t_2)$ existiert, das eine Wettkampfbedingung erzeugt, und $t_1$ und $t_2$ in $m$ aktiviert sind.

\begin{figure}
	\centering
	\begin{tikzpicture}[node distance=2cm,on grid, auto]
		\node[place, label=$p_1$] (p1) {};
		\node[transV, label=$t_1$, right = of p1] (t1) {};
		\node[place, label=$p_2$, tokens=1, above right = of t1] (p2) {};
		\node[place, label=$p_3$, tokens=1, below right = of t1] (p3) {};
		\node[transV,label=$t_2$, right = of p2,label=below:{\LSset{a,{\color{red}b}}{a}} ] (t2){};
		\node[transV,label=$t_3$, right = of p3,label=below:{\LSset{b}{{\color{red}b}}}] (t3){};
		\node[place, label=$p_4$, right = of t2] (p4) {};
		\node[place, label=$p_5$, right = of t3] (p5) {};
		\node[transV, label=$t_4$, below right = of p4] (t4) {};
		\node[place, label=$p_6$, right = of t4] (p6) {};
	
	
	
		\draw 
		(p1) edge[post] (t1)
		(t1) edge[post] (p2)
		(t1) edge[post] (p3)
		(p2) edge[post] (t2)
		(p3) edge[post] (t3)
		(t2) edge[post] (p4)
		(t3) edge[post] (p5)
		(p4) edge[post] (t4)
		(p5) edge[post] (t4)
		(t4) edge[post] (p6)
		;
	\end{tikzpicture}
	\caption[Ein erweitertes Petri-Netz, dessen Markierung eine Wettkampfbedingung enthält.]{Ein erweitertes Petri-Netz, dessen Markierung eine Wettkampfbedingung enthält. Die Schreibmenge von $t_3$ enthält die Variable $b$, diese ist allerdings auch in der Lesemenge von $t_2$ enthalten. Die problematische Variable {\color{red}$b$} ist rot markiert. Das Auftauchen von $b$ in der Lesemenge von $t_3$ stellt kein Problem dar.}\label{fig:wettkampfpetri}
\end{figure}

Um das Konzept zu veranschaulichen, ist in Abbildung~\ref{fig:wettkampfpetri} ein erweitertes Petri-Netz gezeigt, dessen Markierung eine Wettkampfbedingung enthält. Die Plätze $p_2$ und $p_3$ enthalten je eine Markierung. Daher müssen nun die Transitionen betrachtet werden, die aktiviert sind. Diese sind $\{t_2,t_3\}$. Wie in der Abbildung rot markiert, ist ersichtlich, dass sich die Lese- und Schreibmenge der Transitionen überschneiden. Dadurch gilt $V_\text{krit} = \{b\} \neq \emptyset$ und die gezeigte Markierung enthält eine Wettkampfbedingung.

Ein erweitertes Petri-Netz enthält eine Wettkampfbedingung, wenn mindestens eine Markierung seines Erreichbarkeitsgraphen eine Wettkampfbedingung enthält.

Aufgabe der nebenläufigen Programmierung ist es also, \glspl{Anweisung} so zu strukturieren, dass alle von nebenläufigen \glspl{Anweisung} genutzten Daten threadsicher sind oder Daten, die nicht threadsicher sind, explizit synchronisiert werden. Über die Modellierung von Petri-Netzen ausgedrückt, gilt es also die Struktur eines \glsuseri{Programm} so zu definieren, dass in dem Petri-Netz, welches das \gls{Programm} modelliert, keine Wettkampfbedingungen existieren.

\subsubsection{Synchronisierung}
Um das Auftreten von Wettkampfbedingungen beim schreibenden Zugriff von nebenläufigen \glspl{Anweisung} auf geteilte Ressourcen zu vermeiden, muss sichergestellt werden, dass der Endzustand nach der Ausführung der \glspl{Anweisung} unabhängig von der Ausführreihenfolge der daraus resultierenden \glspl{Aktivitaet} ist. Dieser Vorgang wird \emph{Synchronisierung} genannt~\cite[S.~4]{Maurer2019}.

\paragraph{Synchronisierungsanforderungen} Nach \textcite[S.~132~ff.]{Herrtwich1989} gibt es allgemein zwei Arten von Anforderungen an Synchronisierungsmechanismen, je nachdem welche Art von Zugriff auf geteilte Ressourcen stattfindet. Sie unterscheiden dabei zwischen \emph{kausal abhängigen} Relationen und \emph{kausal unabhängigen} Relationen zwischen nebenläufigen \glspl{Anweisung}. Ist \gls{Anweisung} $B$ kausal abhängig von \gls{Anweisung} $A$, so dürfen die \glspl{Aktivitaet} von $B$ erst ausgeführt werden, nachdem die \glspl{Aktivitaet} von $A$ abgeschlossen sind. Beispielsweise werden in der Blocklib Chunks im Hintergrund geladen. Benötigt eine \gls{Anweisung} eines Spielelements einen noch nicht geladenen Chunk, so muss die Ausführung der zugehörigen \gls{Aktivitaet} solange aufgeschoben werden, bis der Chunk geladen ist. Bei kausal unabhängigen Relationen muss sichergestellt werden, dass der Zugriff auf die Ressourcen nicht zur gleichen Zeit stattfindet.  Beispielsweise wird die Liste der zu zeichnenden Objekte in der Blocklib von verschiedenen Systemen angepasst. Hier muss also darauf geachtet werden, dass diese Anpassungen zeitlich getrennt stattfinden. Sonst könnten beispielsweise hinzugefügte Elemente verloren gehen könnten, wenn zwei Threads parallel Elemente hinzufügen, weil einer die Änderungen des anderen überschreibt.

\paragraph{Synchronisierungsmethoden} Es gibt mehrere Möglichkeiten, die Synchronisierung zwischen nebenläufigen \glspl{Anweisung} zu erreichen. \textcite{Michael1996} unterteilen Algorithmen diesbezüglich in zwei Kategorien: \emph{blockierende} und \emph{nicht-blockierende}. Blockierende Synchronisierungsalgorithmen definieren sogenannte \emph{kritische Bereiche}. Ein kritischer Bereich ist eine Menge von \glspl{Anweisung} von deren \glspl{Aktivitaet} nur eine gleichzeitig ausgeführt werden darf. Sind mehrere auszuführende \glspl{Aktivitaet} Teil eines kritischen Bereichs, muss durch \emph{Sperren} (engl. Locks) sichergestellt werden, dass nur eine der \glspl{Aktivitaet} gleichzeitig ausgeführt wird. Möchte ein \gls{Rechenprozess} einen kritischen Bereich betreten, muss er sich die Sperre zunächst aneignen. Ist die Sperre bereits im Besitz eines anderen \glsuseri{Rechenprozess}, muss der \gls{Rechenprozess} warten, bis seine Aneignung der Sperre erfolgreich ist. Dadurch führen solche Synchronisierungsmethoden laut \textcite{Michael1996} möglicherweise dazu, dass \glspl{Rechenprozess} durch andere \glspl{Rechenprozess} beliebig lang beim Zugriff auf geteilte Daten aufgehalten werden können.

Nicht-blockierende Algorithmen dagegen garantieren, dass mindestens einer der beteiligten \glspl{Rechenprozess} mit einer endlichen Anzahl an Schritten endet. Zudem konnte \textcite{Herlihy1991} bereits 1991 zeigen, dass beispielsweise unter Verwendung von atomaren \emph{Compare-and-Swap}-Instruktionen Datenstrukturen implementiert werden können, auf die nicht-blockierend von einer unbegrenzten Zahl an nebenläufigen \glspl{Aktivitaet} zugegriffen werden kann. Durch die Atomarität der genutzten \glspl{Anweisung} wird sichergestellt, dass die geteilten Ressourcen nicht in einen unvorhergesehenen Zustand übergehen. Die eben genannte Instruktion Compare-and-Swap ist eine \gls{Anweisung}, in der atomar der Wert einer Variablen ausgelesen und mit einem übergebenen Wert verglichen wird (Compare). Danach wird, weiterhin atomar, der Inhalt der Variablen mit einem zweiten übergebenen Wert genau dann überschrieben (Swap), wenn die verglichenen Werte identisch waren. Die Instruktion gibt dann den vorherigen Wert zurück. Eine andere Variante der Instruktion (\emph{Compare-and-Set}) gibt zurück, ob der Wert überschrieben wurde oder nicht. Die Atomarität dieser gesamten Instruktion muss durch die Hardware sichergestellt werden, das ist softwareseitig nicht möglich. Alle gängigen \acp{cpu} stellen Instruktionen dieser Art zur Verfügung.

\section{Multithreading in Java}
In Java können Threads über ein einheitliches Interface verwaltet werden. Zur Entkopplung vom Betriebssystem übernimmt die Java Runtime die Aufrufe der betriebssystemspezifischen Funktionen zur Erzeugung und Verwaltung eines Threads~\cite[S.~3]{Friesen2015}. Das einheitliche Interface wird durch die Klasse \class{Thread} repräsentiert, die einen Java Thread von einem Thread des  Betriebssystems entkoppelt. 

Um die Ausführung eines Java Threads zu starten, wird dessen Methode \code{start(): void} aufgerufen~\cite[S.~8]{Friesen2015}. Ohne eine Angabe der auszuführenden Anweisungen ist allerdings noch nicht definiert, welche Aktivitäten der Thread bei der Ausführung durchführen soll. Java enthält das Interface \class{Runnable}, das eine einzige Methode \code{run(): void} definiert, die weder Parameter noch Rückgabewert besitzt. Die Klasse \class{Thread} besitzt Konstruktoren die als Argument ein Objekt erwarten, das \class{Runnable} implementiert, wird der Java Thread gestartet, so wird in der Ausführung die \code{run()} Methode des übergebenen Objekts aufgerufen~\cite[S.~3]{Friesen2015}. Ein einfaches Programm, das auf diese Weise arbeitet, könnte wie folgt aussehen.
\begin{lstlisting}
public static void main(String[] args) {
	var helloThread = new Thread(new Runnable() {
		@Override
		public void run() {
			System.out.println("Hello thread!");
		}
	});
	helloThread.start();
}
\end{lstlisting}
Das Programm erzeugt ein Java Thread Objekt, dem eine Implementierung von \class{Runnable} übergeben wird. In diesem Beispiel wird nur der Text \enquote{Hello thread!} ausgegeben. Hier könnte sich natürliche ein beliebiger Algorithmus befinden, der dann in dem Thread ausgeführt wird.

Alternativ kann eine Klasse von \class{Thread} erben, da \class{Thread} selbst \class{Runnable} implementiert, und die Definition von \code{run()} überschreiben~\cite[S.~335]{Rauber2006}. Da die Klasse von \class{Thread} erbt, besitzen Objekte der Klasse ebenfalls die Methode \code{start()}. Wird diese aufgerufen, so wird wie bei \class{Thread} ein Betriebssystem Thread gestartet, der nun die überschriebene Definition von \code{run()} ausführt. Das Erben von \class{Thread} ist im Normalfall nicht empfohlen, da dieses Design einige Probleme verursacht. Beispielsweise kann die Thread-Unterklasse von keiner anderen Klasse mehr erben~\cite[S.~335]{Rauber2006}, zudem sind die Definition der Anweisungen mit der Definition der Ausführung höchstmöglich gekoppelt, was dem grundlegenden Prinzip der \emph{losen Kopplung} entgegensteht. Entscheidet man sich also für dieses Design, sollte dies eine gute Begründung haben.

Um die Ausführung mehrerer Java Threads zu koordinieren, bietet Java einige Synchronisationsmechanismen. Diese werden nun vorgestellt

\subsubsection{Synchronisierung in Java}
Da die Bandbreite der Möglichkeiten zur Synchronisierung in Java ist sehr groß. Daher werden hier nur einige ausgewählte Methoden dargestellt. Für eine ausführlichere Lektüre zu dem Thema lassen sich die Bücher von \textcite{Friesen2015} und \textcite{Hettel2016} empfehlen.

\paragraph{\texttt{synchronized}} Das Schlüsselwort \code{synchronized} ermöglicht die blockierende Synchronisierung beim Zugriff auf beliebige Objekte. Das Schlüsselwort kann auf zwei Arten eingesetzt werden~\Cite[S.~339~ff.]{Rauber2006}. Einerseits kann ein Code-Block definiert werden in dem der Zugriff auf ein Objekt synchronisiert wird, wie in folgendem Beispiel.
\begin{lstlisting}
public class MultithreadedIncrement {
	public static Integer count = 0;
	public static void main(String[] args) {
		var increment = new Runnable(){
			@Override
			public void run() {
				synchronized (count){
					count++;
					System.out.println(count);
				}
			}
		};
		new Thread(increment).start();
		new Thread(increment).start();
	}
}	
\end{lstlisting}
Das Programm definiert ein statisches Feld \code{count}, dann wird ein \class{Runnable} implementiert, das einen \code{synchronized} Block enthält. Versucht ein Thread, diesen Block auszuführen, versucht er automatisch die zu \code{count} gehörige Sperre erhalten. Dieser Vorgang ist threadsicher und nur ein Thread kann die Sperre zu einem Objekt gleichzeitig besitzen. Dadurch ist es nun möglich, gefahrlos den Wert von \code{count} zu erhöhen und den neuen Wert auf der Konsole auszugeben. Beim Verlassen des \code{synchronized} Blocks wird die Sperre des zugehörigen Objekts wieder freigegeben. Ist die Sperre nicht belegt, wenn ein Thread diese erhalten möchte, wird die Ausführung des Threads solange suspendiert, bis diese freigegeben wird. Java ordnet jedem Objekt implizit eine solche Sperre zu~\cite{Friesen2015}, sodass \code{synchronized} mit beliebigen Objekten genutzt werden kann. Für primitive Datentypen ist das nicht möglich.

Das Schlüsselwort \code{synchronized} kann auch als Teil der Signatur einer Methode verwendet werden, in diesem Fall wird kein Objekt angegeben, dessen Zugriff synchronisiert werden soll. Bei Instanzmethoden ist das automatisch das Objekt selbst, das die Methode ausführt also \code{this}, bei statischen Methoden wird der Zugriff auf das Klassenobjekt synchronisiert.

\paragraph{\texttt{volatile}} Variablen können mit dem Schlüsselwort \code{volatile} ausgestattet werden. Threads, die auf eine als \code{volatile} gekennzeichnete Variable zugreifen, lesen immer den aktuellen Wert dieser Variable und greifen dafür immer auf den Hauptspeicher zu anstatt auf den Cache (der schon veraltet sein könnte)~\cite[S.~30~ff.]{Friesen2015}. Lese- und Schreibvorgänge einer \code{volatile} Variablen sind atomar, und damit threadsicher.

Dies gilt insbesondere auch für 64-bit Variablen. Obwohl einige Quellen davor warnen \code{volatile} für Variablen des Typs \code{long} oder \code{double} zu verwenden (unter dem Verweis darauf, dass 32-bit Computer zwei Operationen für das Lesen und Schreiben dieser Variablen benötigen~\cite[S.~34]{Friesen2015}), beschreibt die Java Spezifikation explizit, dass auch der Zugriff auf 64-bit Variablen immer atomar ist, sofern diese als \code{volatile} deklariert sind~\cite{Java7Spec17}. Weiter wird dort insbesondere darauf hingewiesen, dass Entwickler aus diesem Grund 64-bit Variablen, auf die von mehreren Threads zugegriffen wird, als \code{volatile} zu deklarieren, um Atomarität von Lese- und Schreibzugriffen zu gewährleisten. Lesen und subsequentes schreiben einer \code{volatile} Variablen (beispielsweise \code{count++;}) ist allerdings weiterhin nicht atomar, da mehr als nur eine Operation durchgeführt wird -- Lesen, neuen Wert berechnen und Schreiben. \code{volatile} ist also nur dann zu verwenden, wenn der neue Wert einer Variablen nicht vom alten Wert abhängt, beispielsweise, wenn er gänzlich neu berechnet wird. 

\paragraph{Atomics und Concurrent Collections} Da das Ändern von Variablen basierend auf ihrem vorherigen Wert häufig nötig ist (beispielsweise für die Generierung einer fortlaufenden Nummer) und die ständige Synchronisierung der Variable über \code{synchronized} oder andere blockierende Methoden beträchtliche Performanceeinbußen zur Folge haben können~\cite[S.~130]{Friesen2015}, enthält Java \emph{Atomics}. Atomics sind Wrapper für bestimmte Datentypen, die es erlauben, atomare Operationen auch auf Basis des aktuellen Werts einer Variablen durchzuführen ohne dabei den Zugriff auf diese Variable durch eine Sperre zu blockieren~\cite[S.~130]{Friesen2015}. 

Java stellt unter anderem die Klassen \class{AtomicBoolean}, \class{AtomicInteger}, \class{AtomicLong} und \class{AtomicReference}
bereit~\cite[S.~131]{Friesen2015}. Klassen für die Datentypen \code{byte}, \code{short}, \code{float} und \code{double} gibt es nicht. Die Java Dokumentation erklärt dies damit, dass diese über \class{AtomicInteger} und im Fall von \code{double} über \class{AtomicLong} dargestellt werden können~\cite{Java7DocAtomic}.

Eine nützliche Anwendung von Atomics ist die threadsichere Beschränkung einer Klasse auf eine Objektinstanz, ohne globalen Zugriff auf diese bereitstellen zu müssen wie es bei der Standardimplementierung eines \gls{singleton} der Fall wäre. Allerdings ermöglicht diese Implementierung keine Prüfung während des Kompiliervorgangs, sondern erst zur Laufzeit.
\begin{lstlisting}
public class LocalSingleton {
private static AtomicBoolean instanced = new AtomicBoolean(false);

public LocalSingleton() {
	ensureUniqueInstance();
}

private void ensureUniqueInstance(){
	if(!instanced.compareAndSet(false, true)){
		throw new UnsupportedOperationException();
	}
}
\end{lstlisting}
Die gezeigte Klasse \class{LocalSingleton} nutzt eine statische Instanz von \class{AtomicBoolean}, um im Konstruktor zu überprüfen, ob bereits eine Instanz von \class{LocalSingleton} erzeugt wurde. Ist das der Fall, wird eine Exception geworfen und das Objekt kann nicht erzeugt werden. Dazu wird eine Compare-and-Swap Instruktion verwendet, die in Java \code{compareAndSet(...)} genannt wird.

Compare-and-Swap Instruktionen werden ebenfalls in den von Java bereitgestellten \emph{Concurrent Collections} verwendet~\cite[S. 133~f.]{Friesen2015}. Es gibt verschiedene nicht-blockierende Implementierungen für Warteschlangen (\class{ConcurrentLinkedQueue}), Maps (\class{ConcurrentHashMap}), Mengen (\class{ConcurrentSkipListSet}) und weiteres. Diese sind selbstverständlich besonders nützlich, wenn verschiedene Threads auf solche Datenstrukturen zugreifen müssen, da keine explizite Synchronisierung nötig ist und keine Sperren verwendet werden, welche sich zu einem Flaschenhals entwickeln könnten~\cite[S.~125~ff.]{Friesen2015}.

\subsubsection{Executors}\label{sec:executor} Neben Möglichkeiten zur Synchronisierung bietet das Concurrency Framework in Java zudem Klassen und Interfaces, die die Verwaltung und Ausführung nebenläufiger Aufgaben vereinfachen. Die Interfaces \class{Executor} und \class{ExecutorService} liefern eine einheitliche Schnittstelle, um die Übergabe von nebenläufigen Anweisungen von der Verwaltung der tatsächlichen Ausführung der Anweisungen zu trennen~\cite[S.~70~ff.]{Friesen2015}. Das Interface \class{Executor} definiert genau eine Methode \code{execute(command: Runnable): void}. Die Klasse, die einen Executor nutzt, hat keine Kenntnis über die Anzahl und Beschaffenheit der Threads, die die Anweisungen tatsächlich ausführen. Tatsächlich ist nicht einmal bekannt, ob die Anweisungen überhaupt in Threads ausgeführt werden. Nutzt ein Executor im Hintergrund Threads für die Ausführung der Anweisungen, so werden diese auch \emph{Worker Threads} genannt.

Das Interface \class{Runnable} definiert eine Schnittstelle für die Definition nebenläufiger Anweisungen, das Konzept der asynchronen Ausführung von Anweisungen, die ein Ergebnis liefern, wird davon jedoch nicht abgedeckt. Diese Lücke schließt das Interface \class{Callable<V>}, das die generische Methode \code{call(): V} enthält und somit beliebige Rückgabewerte ermöglicht. \textcite{Friesen2015} nennt ein Objekt, das entweder das Interface \class{Runnable} oder das Interface \class{Callable} implementiert, \emph{Task}. Diese Konvention wird hier übernommen. 

Analog zur Erweiterung der Möglichkeiten durch \class{Callable} erweitert das Interface \class{ExecutorService} die Möglichkeiten der Ausführung von Tasks. Während das Interface \class{Executor} lediglich die Ausführung von Runnables ermöglicht, kann ein Objekt, das \class{ExecutorService} implementiert, auch die Ausführung von Callables verwalten. Zudem definiert das Interface Methoden, die den Lebenszyklus der implementierenden Objekte selbst betreffen, wie beispielsweise \code{shutdown(): void}. Bei der Übergabe eines Callables an einen \class{ExecutorService} gibt dieser ein Objekt \class{Future<V>} zurück, das die Interaktion mit der asynchronen Anweisung ermöglicht. Ein vollständiges Klassendiagramm der beschriebenen Interfaces ist in Anhang~\vref{appendix:concFrameworkExecutor} enthalten.

\subsubsection{Koordinierung asynchroner Tasks}\label{sec:CompletableFuture}

Die Methoden von \class{Future<V>} ermöglichen zwar gewisse Interaktionen, wie mittels \code{get(): V} auf die Fertigstellung der asynchronen Anweisung zu warten und deren Rückgabewert zu erhalten, Möglichkeiten zur Komposition voneinander abhängiger Tasks bietet das Interface allerdings nicht. Es ist, wie \textcite[S.~240]{Hettel2016} es nennt, ein \emph{pullbasiertes} Interface. Ein Entwickler muss also explizit abfragen, ob ein asynchroner Task bereits abgeschlossen ist. Dies erzeugt die Gefahr Threads zu blockieren, da sie auf die Werte und damit auf die Fertigstellung anderer Tasks warten müssen~\cite[S.~239]{Hettel2016}.

Die Klasse \class{CompletableFuture<V>} implement dagegen zusätzlich das \emph{pushbasierte} Interface \class{CompletionStage}. Sie übernimmt also die Verwaltung der Komposition asynchroner Tasks~\cite[S.~240~ff.]{Hettel2016}. Damit wird die Verkettung, Aufspaltung, Zusammenführung und Auswahl von Tasks ermöglicht~\cite[S.~250~ff.]{Hettel2016}. Mit Verkettung ist die Ausführung eines Tasks nach einem anderen gemeint. Aufspaltung bezeichnet die Ausführung mehrerer Tasks nach der Vollendung eines Tasks. Zusammenführung ist die Ausführung eines Tasks, nachdem mehrere andere Tasks abgeschlossen sind. Mit Auswahl ist die Möglichkeit gemeint, einen Task zu starten, nachdem der schnellere von zwei vorherigen Tasks abgeschlossen ist. Ein Klassendiagramm von \class{CompletionStage} ist in Anhang~\vref{appendix:CompletionStage} zu finden. Dort den genannten Kompositionsmöglichkeiten die jeweiligen Methoden zugeordnet.
\section{Multithreading in Spielen}
Computerspiele nutzen Multithreading seit dem Aufkommen von Mehrkern-Prozessoren, um eine bessere Perforance zu erreichen~\cite{Davies2006}. Es gibt verschiedene für Spiele geeignete Multithreading-Architekturen, die jeweils eigene Vor- und Nachteile mit sich bringen. Im folgendenden werden zwei weit verbreitete Multithreading Architekturen betrachtet.

\subsection{Thread pro System}

\begin{figure}
	\centering
	\begin{tikzpicture}
			\fill[lightgray]  (0,0) rectangle (11,1);
			\fill[lightgray] (0,-1.5) rectangle (11,-0.5);
			\fill[lightgray]  (0,1.5) rectangle (11,2.5);
			
			\node at (-1,2) {Thread 1};
			\node at (-1,0.5) {Thread 2};
			\node at (-1,-1) {Thread 3};
			
			
			\fill [orange] (0.5,2.5) rectangle (4.5,1.5);
			\fill [orange] (5,2.5) rectangle (9,1.5);
			\fill [magenta] (0.5,1) rectangle (3.5,0);
			\fill [magenta] (5,1) rectangle (8,0);
			
			
			\node at (2.5,2) {Simulation $n$};
			\node at (7,2) {Simulation $n+1$};
			\node at (2.5,3.5) {Frame $n$};
			\node at (7,3.5) {Frame $n+1$};
			
			\draw  (0.5,3) rectangle (4.5,-2);
			\draw  (5,3) rectangle (9,-2);
			
			\node at (2,0.5) {Render $n-1$};
			\node at (6.5,0.5) {Render $n$};
			\node at (4,0.5) {?};
			\node at (8.5,0.5) {?};
			\node at (2.5,-1) {?};
			\node at (7,-1) {?};
	\end{tikzpicture}
	\caption{Systeme laufen auf eigenen Threads. Es wird die Auslastung während zwei Frames $n$ und $n+$ gezeigt. Einige Threads sind dadurch nicht zu \SI{100}{\percent} ausgelastet, was durch Fragezeichen (?) gekennzeichnet ist. In diesem Fall kann Thread 3 gar nicht genutzt werden. Thread 2 übernimmt das Rendering, durch Pipelining wird dabei immer der Zustand der Simulation des letzten Frames dargestellt. Die Abbildung ist von einer Grafik aus \cite[S.~14]{Tatarchuk2014} inspiriert.}\label{fig:sot}
\end{figure}
Eine einfache Architektur vergibt, wie in Abbildung~\ref{fig:sot} gezeigt, für einzelne Systeme (wie Simulation und Rendering) eigene Threads, die diese exklusiv nutzen~\cite{Davies2006,Tatarchuk2014,Genova2015,Hodgman2016}. Durch Pipelining kann das Rendering System dann beispielsweise auf den Spielzustand des vorherigen Frames für die Darstellung zugreifen. Diese Architektur bietet einige Vor- und Nachteile, die nun erörtert werden.
\begin{itemize}
	\item[$+$] Die Architektur bietet den Vorteil, dass die Ausführung innerhalb der Systeme seriell ist, was die Entwicklung wie in Abschnitt~\ref{sec:nebenl-folgen} beschrieben einfacher macht, als nebenläufige Entwicklung.
	\item[$+$] Solange die Systeme voneinander separiert sind, gibt es keine Probleme die mit Nebenläufigkeit einhergehen, da beispielsweise Wettkampfbedingungen nur auftreten können, wenn nebenläufig auf geteilte Daten zugegriffen wird. Da keine Wettkampfbedingungen verhindert werden müssen, entfällt (großteils) und Deadlocks sind ausgeschlossen.
	\item[$-$] Auf heterogenen Plattformen, die beispielsweise unterschiedliche Prozessoreinheiten besitzen, kann es zu Performance Problemen kommen. Die Architektur muss eventuell für jede Plattform anders aufgebaut sein. Es gibt also eine gewisse Abhängigkeit von der Hardware.
	\item[$-$] Synchronisierung der Systeme erfordert einen großen Speicheraufwand, da alle Daten (beispielsweise mit einem Double Buffer) zwischengespeichert werden müssen.
\end{itemize}

Die Nutzung von Threads pro System bietet sich an, wenn nur für eine Plattform, beispielsweise PC, entwickelt wird und es keine starken Speicherplatzrestriktionen gibt. Zudem ist die Implementierung einfacher als die der folgenden Architektur.

\subsection{Jobsystem}

\begin{figure}
	\centering
	\begin{tikzpicture}
		\fill[lightgray]  (0,0) rectangle (11,1);
		\fill[lightgray] (0,-1.5) rectangle (11,-0.5);
		\fill[lightgray]  (0,1.5) rectangle (11,2.5);
		
		\node at (-1,2) {Thread 1};
		\node at (-1,0.5) {Thread 2};
		\node at (-1,-1) {Thread 3};
	
		\foreach \i in {0,3.5,7}{
		\fill [orange,draw=lightgray] ($(\i,0) + (0.5, 1.5)$) rectangle ($(\i,0) +(1.5, 2.5)$);
		\fill [orange,draw=lightgray] ($(\i,0) + (0.5,-1.5)$) rectangle ($(\i,0) +(1.5,-0.5)$);
		\fill [orange,draw=lightgray] ($(\i,0) + (1.5, 1.5)$) rectangle ($(\i,0) +(2.5, 2.5)$);
		\fill [orange,draw=lightgray] ($(\i,0) + (0.5, 0.0)$) rectangle ($(\i,0) +(1.5, 1.0)$);
		
		\fill [magenta,draw=lightgray] ($(\i,0) + (2.5, 1.5)$) rectangle ($(\i,0) + (3.5, 2.5)$);
		\fill [magenta,draw=lightgray] ($(\i,0) + (2.5, 0.0)$) rectangle ($(\i,0) + (3.5, 1.0)$);
		\fill [magenta,draw=lightgray] ($(\i,0) + (2.5,-1.5)$) rectangle ($(\i,0) + (3.5,-0.5)$);
		
		\draw  ($(\i,0) + (0.5,3)$) rectangle ($(\i,0) + (3.5,-2)$);
		
		\node[font=\footnotesize] at ($(\i,0) + (1,2)$) {Sim 1};
		\node[font=\footnotesize] at ($(\i,0) + (1,0.5)$) {Sim 2};
		\node[font=\footnotesize] at ($(\i,0) + (1,-1)$) {Sim 3};
		\node[font=\footnotesize] at ($(\i,0) + (2,2)$) {Sim 4};
		\node[font=\footnotesize] at ($(\i,0) + (3,2)$) {Ren 1};
		\node[font=\footnotesize] at ($(\i,0) + (3,0.5)$) {Ren 2};
		\node[font=\footnotesize] at ($(\i,0) + (3,-1)$) {Ren 3};
	
		}
	\node at (2,3.5) {Frame $n$};
	\node at (5.5,3.5) {Frame $n+1$};
	\node at (9,3.5) {Frame $n+2$};
	\end{tikzpicture}
	\caption{Die Systeme definieren Jobs, die frei auf die vorhandenen Threads aufgeteilt werden. Sim steht für Simulation und Ren für Rendering. Die Abbildung ist schematisch, die Jobs müssen nicht exakt die gleiche Länge haben, auch wenn kurze Jobs mit ähnlichen Laufzeiten ideal sind. Im Vergleich zu Abbildung~\ref{fig:sot} lässt sich eine deutlich höhere Auslastung der Threads erkennen, was in höheren FPS zeigt.}\label{fig:jobt}
\end{figure}
\todo{Abbildung Caption}
Eine etwas komplexere Architektur liefert die Implementierung eines sogenannten \emph{Job-} oder \emph{Task-Systems}~\cite{Davies2006,Tatarchuk2014,Genova2015,Hodgman2016}. Ein Schema der Architektur ist in Abbildung~\ref{fig:jobt} zu sehen. Anstatt Aufgaben fest auf einzelne Threads zu verteilen, gibt es \emph{Jobs} und ein System, das die Jobs auf Threads verteilt. Konzeptuell entspricht die Architektur sehr dem in Abschnitt~\ref{sec:executor} beschriebenen \class{Executor} in Java. 

Die idealerweise kurzen Jobs werden an das System übergeben. Das System verwaltet die Threads, deren Anzahl idealerweise der Anzahl der Hardwarethreads entspricht, da dann die Leistung des Prozessors optimal ausgenutzt werden kann und die Anzahl der Kontextwechsel minimiert wird. Ist ein Thread im Leerlauf, wird ein Job an den Thread zur Ausführung übergeben.

\begin{figure}
	\centering
	\begin{tikzpicture}
		\node[fill=orange] (sim1) {Sim 1};
		\node[fill=orange,below=of sim1] (sim2)  {Sim 2};
		\node[fill=orange,below=of sim2] (sim3)  {Sim 3};

		\node[fill=orange,right=of sim2] (sim4)  {Sim 4};

		\node[fill=magenta,right=of sim4] (ren2)  {Ren 2};
		\node[fill=magenta,above=of ren2] (ren1)  {Ren 1};
		\node[fill=magenta,below=of ren2] (ren3)  {Ren 3};

		
		\draw[->](sim1.east) -- (sim4.north west);
		\draw[->](sim2.east) -- (sim4.west);
		\draw[->](sim3.east) -- (sim4.south west);
		\draw[->](sim4.north east) -- (ren1.west);
		\draw[->](sim4.east) -- (ren2.west);
		\draw[->](sim4.south east) -- (ren3.west);
	\end{tikzpicture}
	\caption{Beispiel eines Job Graphen, der der Ausführung in Abbildung~\ref{fig:jobt} zugrundeliegen könnte. Die farbigen Knoten stellen Jobs dar, wobei orange Knoten Simulationsjobs und magentafarbene Knoten Renderingjobs sind. Die gerichteten Kanten zeigen die mögliche Ausführungsreihen folge an. So muss beispielsweise Sim 4 nach Sim 3 ausgeführt werden. Gibt es keinen Pfad von Job $A$ zu Job $B$ so können beide parallel ausgeführt werden.}\label{fig:jobdependencies}
\end{figure}
\todo{Abbildung}
Das System kann zudem so implementiert werden, dass es Abhängigkeiten zwischen einzelnen Jobs gibt, das heißt, dass ein Job beispielsweise erst ausgeführt wird, wenn ein anderer Job abgeschlossen ist. Daraus bildet sich ein Job Graph, der die nebenläufige Ausführung der Jobs beschreibt. Es hat sich als hilfreich erwiesen, feste Synchronisationspunkte zu definieren, damit auf systemübergreifende Daten zugegriffen werden kann. Das kann im Jobsystem wie in Abbildung~\ref{fig:jobdependencies} durch Jobs repräsentiert werden, die Abhängigkeiten zu vielen vorangehenden Jobs sammeln und dann als Abhängigkeit für die folgenden Jobs dienen. Betrachten wir nun einige Vor- und Nachteile eines Jobsystems.
\begin{itemize}
	\item[$+$]  Da die Durchführung von Jobs nicht an bestimmte Threads gekoppelt ist, kann Hardware mit mehr Prozessoreinheiten sofort voll ausgenutzt werden. Wenn die Anzahl der Prozessoren zur Laufzeit ermittelt wird, funktioniert das sogar dynamisch. Ist eine Prozessoreinheite vorhanden, wird diese genutzt, gibt es $100$ Prozessorkerne, werden $100$ genutzt. Durch ein Jobsystem kann also jede beliebige in Hardware vorhandene Parallelität ausgenutzt werden. Dadurch werden ohne zusätzlichen Programmieraufwand verschiedenste Plattformen unterstützt.
	\item[$+$] Ein Jobsystem bietet die Möglichkeit, Simulation und Rendering in einem Frame sequentiell durchzuführen, ohne alle Nebenläufigkeit zu verlieren. Durch das Jobsystem können die Systeme intern parallelisiert werden. Werden Simulation und Rendering sequentialisiert, ist es nicht nötig des Spielzustand zwischenzuspeichern, da die Simulation während des Renderings bereits abgeschlossen ist und somit Wettkampfbedingungen ausgeschlossen sind (da das Rendering nur lesend auf Daten zugreift). Damit benötigt man weniger Speicher als bei einer Architektur, die auf einen Zwischenspeicher angewiesen ist. Spiele, die eh schon viel Speicher verbrauchen und auf speicherbeschränkten Systemen ausgeführt werden sollen, können auf diese Art von Nebenläufigkeit guten Gebrauch machen.
	\item[$+$] Wird ein neues System hinzugefügt, beispielsweise ein System zur Verwaltung der Audioausgabe, muss die grundlegende Architektur nicht angepasst werden, um das neue System nebenläufig in das Spiel zu integrieren. Das System kann einfach Jobs erzeugen und diese zur Ausführung bringen.
	\item[$-$] Das Jobsystem basiert auf Jobs. Diese müssen erstellt werden. Soll also ein bestehendes Spiel ein Jobsystem nutzen, muss der gesamte Code so umgeschrieben werden, das bestehende Abläufe in kleine Jobs aufgeteilt werden und dann mit den korrekten Abhängigkeiten nebenläufig ausgeführt werden. Es muss also tief in die bestehenden Systeme eingegriffen werden. Dieses Problem ergibt sich bei der \ac{sot} Architektur nicht.
\end{itemize}

\chapter{Analyse der Blocklib}

\section{Game Loop}
Ein nahezu universelles Designelement von Spielen ist die sogenannte \emph{Game Loop}~\cite{}. Sie ist für die Aktualisierung des Spiels zuständig und umfasst typischerweise die Aufgaben 
\begin{itemize}
\item Verarbeitung der Eingaben,
\item Berechnung des neuen Spielzustandes und
\item (grafische) Ausgabe.
\end{itemize}

\section{Nutzung von Multithreading}
Bereits vor dieser Arbeit werden in der Blocklib Threads genutzt, um nebenläufige Anweisungen zu parallelisieren. Darunter fallen insbesondere Hintergrundaufgaben, wie das Erstellen von Chunks, und das Ausführen zeitlich gesteuerter Abläufe, wie die Updates der Flüsse von Flüssigkeiten. Insgesamt gibt es zu Beginn der Arbeit sechs Klassen, die Threads erzeugen und nutzen. Eine Auflistung der Klassen und deren Verwendung von Threads ist in Tabelle~\ref{tab:concTasksBlocklib} zu finden.

\begin{table}
	\renewcommand{\arraystretch}{1.5}
	\begin{tabularx}{\textwidth}{lX}
		\toprule
		Klasse & Bestehende Nutzung von Nebenläufigkeit \\
		\midrule
		\class{ChunkStorage} & Executors werden genutzt, um die Generierung von Chunks und das Meshing von Chunks auszulagern.\\
		\class{EventManager} & Eventhandling wird in Threads ausgelagert, um Deadlocks zu verhindern, wenn die benachrichtigten Beobachter versuchen, auf die selben kritischen Bereiche zuzugreifen. \\
		\class{ConnectionInfo} & Im Server wird ein wiederkehrender Task gestartet, um Ping Nachrichten an die Clients zu senden.\\
		\class{NioClientNetworkLayer} & Ein Thread, der im Hintergrund mit dem Server kommuniziert, wird gestartet.\\
		\class{NioServerNetworkLayer} & Ein Thread, der im Hintergrund mit den Clients kommuniziert, wird gestartet.\\
		\class{Fluid} & Ein Thread wird gestartet, der den Zustand der Flüssigkeit berechnet.\\
		\bottomrule
	\end{tabularx}
	\caption{Nebenläufige Anweisungen in der Blocklib}\label{tab:concTasksBlocklib}
\end{table}

Es gibt keine Zentrale Stelle an der die Erstellung von Threads verwaltet wird. Jede Klasse kümmert sich selbst darum Threads zu erzeugen und zu beenden. Dadurch ist die Anzahl der genutzen Threads nicht leicht zu überblicken. Dadurch lässt sich zum einen schwierig sicherstellen, dass die Anzahl der genutzten Threads die Anzahl der Prozessoren nicht überschreitet, was zu unnötigen Kontextwechseln führt und somit die Performance verschlechtert, zum anderen können die bestehenden Threads möglicherweise nicht voll ausgelastet werden, da nur einzelne Klassen auf diese Zugriff haben.


\chapter{Entwicklung einer Threading API}


\subsection{Threading in Java}

\section{Anforderungen der blocklib an eine Threading API}
\subsection{Kontrolle über Threads an einer Stelle}
\input{chapters/Anforderungen/Kontrolle.tex}
\subsection{Verteilung von Berechnungen an Worker Threads}
\subsection{Verhinderung von Livenessproblemen durch Archtektur}
\subsection{Einfaches starten von Hintergrundtasks und einmaligen Aufgaben}

\section{Design der Threading API}
Bei der Nutzung von Multithreading in Spielen gibt es zwei prominente Ansätze, zum einen \ac{sot}, in Bezug auf Rendering also die Nutzung eines separaten Threads, der das Rendern übernimmt, und zum anderen den Einsatz einer Jobarchitektur.

Die Idee der Nutzung eines Render-Threads, wie aus \ac{sot} bekannt, ergibt sich daraus, dass Simulation des Spiels und Anzeige zwei unabhängige Bereiche sind, die nebenläufig ausgeführt werden können. Da das Rendering sehr rechenintensiv ist, ist es also sinnvoll diese \glspl{Aktivitaet} in einen eigenen Thread auszulagern, damit sie den gesamten Frame für die Ausführung nutzen können. Eine Jobarchitektur ermöglicht es, neben den \glspl{Aktivitaet} des Renderings auch die Simulation in nebenläufige \glspl{Aktivitaet} aufzuteilen. 

Da es sich bei der Blocklib, mit knapp $60000$ Zeilen Code und über $900$ Java Dateien, um ein großes bereits bestehendes System handelt, ist es unrealistisch, in kurzer Zeit das gesamte \gls{Programm} so umzustrukturieren, dass es vollständig eine Jobarchitektur nutzt. Diese Umstrukturierung erfordert tiefes Verständnis für jedes zu ändernde System. Zudem müssen viele Bereiche grundlegend geändert werden, um die Abhängigkeiten der Systeme zu verringern beziehungsweise auszuschließen.

Die Komplexität der Blocklib wird kurz an einem Beispiel erläutert. Die Blocklib enthält eine Klasse \class{Context}, die als eine Art \gls{singleton}-Implementierung eines \glspl{servicelocator}~\cite[S.~301~ff.]{Nystrom2015} dient. Der \gls{servicelocator} ist ein Entwurfsmuster zur dynamischen Auflösung von Abhängigkeiten über ein zentrales Modul. Da \glspl{singleton} ähnlich wie globale Variablen von überall aufgerufen werden können, wird der Code dadurch schwieriger nachvollziehbar~\cite[S.~108]{Nystrom2015}. Dadurch und durch die ebenfalls resultierende Kopplung unterschiedlicher Komponenten würden bei einer naiven Umsetzung der Jobarchitektur unvorhersehbar Wettkampfbedingungen auftreten.
Denn sobald nebenläufige \glspl{Aktivitaet} auf dieselben Ressourcen zugreifen, können Wettkampfbedingungen entstehen. Durch den \class{Context} ist es nun schwierig, einen Überblick zu haben von wo aus auf welche Ressourcen zugegriffen wird. Die Wahrscheinlichkeit, dass Zugriffe entstehen, die zu Wettkampfbedingungen führen, ist also sehr hoch, solange dabei nicht sorgsam vorgegangen wird.

Die Architektur muss also beachten, dass es zwar ein Jobsystem geben soll, dieses aber nicht ausschließlich genutzt wird. Da der Performancegewinn, in Form von \ac{fps}, bei der Nutzung eines Render-Threads als hoch zu erwarten ist, soll dieser in die Blocklib integriert werden. Um die Anforderungen von Kapitel~\ref{sec:anforderungen} zu erfüllen, wird ein Jobsystem implementiert. Das Design der Architektur bedient sich also sowohl bei dem Konzept \ac{sot} also auch dem Jobsystem.
\subsection{Renderthread und Job-API}
Da es sich bei der Blocklib um ein großes bereits bestehendes System handelt \todo{Zeilenzahl angeben}, ist es unrealistisch, in kurzer Zeit das gesamte Programm so umzustrukturieren, dass es vollständig eine Jobarchitektur nutzt. Diese Umstrukturierung erfordert tiefes Verständnis für jedes zu ändernde System. Zudem müssen viele Bereiche grundlegend geändert werden, um die Abhängigkeiten der Systeme zu veringern beziehungsweise deren Vorhandensein auszuschließen.

Die Blocklib enthält beispielsweise eine Klasse \class{Context}, die als eine Art Singletonimplementierung eines Service Locators~\cite[S.~301~ff.]{Nystrom2015} dient. Da Singletons ähnlich wie globale Variablen von überall aufgerufen werden können, wird der Code dadurch potenziell schwieriger nachvollziehbar~\cite[S.~108]{Nystrom2015}. Dadurch und durch die ebenfalls resultierende Kopplung unterschiedlicher Komponenten würden bei einer naiven Umsetzung der Jobarchitektur unvorhersehbar Wettkampfbedingungen auftreten.
Man erinnere sich, dass Wettkampfbedingungen auftreten wenn nebenläufige Aktivitäten auf die selben Ressourcen zugreifen. Durch den \class{Context} ist es nun schwierig, einen Überblick zu haben von wo aus auf welche Ressourcen zugegriffen wird. Die Wahrscheinlichkeit, dass also Zugriffe existieren, die zu Wettkampfbedingungen führen ist also sehr hoch, solange dabei nicht sorgsam vorgegangen wird.

Da der Performancegewinn, in Form von \ac{fps}, bei der Nutzung eines Renderthreads als hoch zu erwarten ist, soll dieser in die Blocklib integriert werden. Um die Anforderungen von Kapitel~\ref{sec:anforderungen} zu Erfüllen, ein Jobsystem implementiert werden, das aber nur an ausgewählten Stellen mit der Blocklib konsolidiert wird und ansonsten für die zukünftige Nutzung bereitsteht.

\subsection{Design des Renderthreads}
Der Renderthread kann einerseits als Teil des Jobsystems designed werden, andererseits gibt es die Möglichkeit den Renderthread von diesem zu trennen. Ist der Renderthread Teil des Jobsystems, kann dieser voll zur Bearbeitung von Jobs mitgenutzt werden. Da üblicherweise die Anzahl der Threads der Anzahl der Prozessorkerne entspricht, kann man so automatisch immer die Leistung aller Prozessorkerne nutzen. Trennt man den Thread dagegen ab, entsteht die Problematik zu entscheiden, wann welche Anzahl von Threads genutzt wird, um möglichst alle Kerne zu nutzen, aber gleichzeitig zu verhindern, dass sich die Threads in der Ausführung gegenseitig behindern. Andererseits gestaltet sich die Implementierung eines getrennten Renderthreads als deutlich einfacher und intuitiver~\cite{Tatarchuk2014}, im Gegensatz zur Jobsystemintegration.

Da das Jobsystem keine vollständig Integration in die Blocklib erfährt und das Rendering in etwa \SI{50}{\percent}\todo{messen} der Framezeit benötigt, bietet es sich an, den Renderthread zu separieren. Der mögliche Performancegewinn durch die Integration in das Jobsystem ist bereits ausgeschlossen, da das Rendering selbst sie gesamte Rechenleistung des Kerns jeden Frame beansprucht. Damit kann auch die Problematik der Anzahl der Threads gelöst werden, indem die Anzahl der Jobthreads um eins veringert wird ohne dadurch Prozessorkerne ohne Arbeit zu verursachen.

Da der Renderthread auf Daten des Spiels zugreift, um die sichtbaren Elemente zu zeichnen, muss sichergestellt werden, dass diese Daten keinen Wettkampfbedingungen unterliegen. In der Spieleentwicklung ist es üblich einen \emph{Spielzustand} (engl. Game State) zu definieren, der während der Simulation in jedem Frame angepasst wird. Der Teil des Zustandes auf den der Renderthread zugreift muss also konstant sein. Während der Simulation wird dieser allerdigs verändert. Einge Möglichkeit, dieses Problem zu beheben, besteht darin einen \emph{Double Buffer}~\cite[S.~143]{Nystrom2015} zu nutzen, um den gesamten Spielzustand zwischenzuspeichern~\cite{Tatarchuk2014}. Die Blocklib ist nicht mit diesem Hintergrund entwickelt worden, weswegen es keine einfache Möglichkeit gibt, den gesamten Spielzustand auf diese Weise zwischenzuspeichern. Die Objekte, die den Spielzustand darstellen sind über die Blocklib hinweg verteilt. 

Um dennoch einen Renderthread nutzen zu können müssen die Objekte identifiziert werden, die für das Rendering benötigt werden. Für diese Objekte muss dann ein geeigneter Double Buffer erzeugt werden, sodass die Daten aus Renderthread-Sicht konstant sind.

\section{Implementierung der Threading API}
Ebenso wie der Entwurf, ist auch die Implementierung der nebenläufigen Architektur prinzipiell in zwei Bereiche unterteilt, die Implementierung des Render-Threads und die Implementierung der Job \acs{api}.

Um einen Render-Thread in die Architektur der Blocklib zu integrieren, muss besonders darauf geachtet werden, dass Wettkampfbedingungen vermieden werden. Wie in Abschnitt~\ref{sec:desgignRender-Thread} beschrieben, kann dies unter anderem durch die Nutzung von Double-Buffern gelingen. Folgend wird die Implementierung des Render-Threads beschrieben. Im darauffolgenden Abschnitt wird die Implementierung des Jobsystems erläutert.
\paragraph{Schattenflackern}
Bei der parallelen Ausführung des Renderthreads ist bei der Darstellung der Schatten ein Flackern zu erkennen. Die Tatsache, dass das Flackern in irregulären Intervallen auftritt, lässt auf das Vorhandensein einer Race-Condition schließen. Diese Vermutung kann dadurch bestätigt werden, dass das Flackern nichtmehr auftritt sobald \code{update()} und \code{render()} sequentialisiert werden.

Da Variablen zeitgleich von verschiedenen Threads bearbeitet werden, ist die Suche nach dem Ursprung des Flackerns schwierig. Da der Fehler nur bei nebenläufiger Ausführung auftritt, müssen die problematischen Stellen auf \code{update()} und \code{render()} verteilt sein. Da der Fehler visuell sichtbar ist, lässt sich das Problem über die Anweisungen in \code{render()} nachverfolgen. 

Eine naive Herangehensweise mittels Debugging ist nahezu unmöglich, da das Flackern zufällig auftritt. Zuerst muss also eine Möglichkeit gefunden werden, um das Auftreten des Flackerns programmatisch zu erkennen.

In der \class{Configuration}-Klasse gibt es das Flag \const{SHOW_SHADOW_MAP_FRAME}. Ist dieses auf \code{true} gesetzt, wird in der linken oberen Ecke des Fensters ein Feld angezeigt, das die sogenannte \class{ShadowMap} zeigt. Abbildung~\ref{fig:ShadowMap} zeigt einen Screenshot der Blocklib mit \code{SHOW_SHADOW_MAP_FRAME = true}. Die \class{ShadowMap} wird genutzt, um zu berechnen, an welchen Stellen Schatten gerendert werden.
\begin{figure}
	\caption{Screenshot}\label{fig:ShadowMap}
\end{figure}
\textcite{Ebbinger2018} beschreibt, wie genau das Rendering von Schatten in der Blocklib umgesetzt ist. Wie der Name der Klasse \class{ShadowMap} vermuten lässt, wird für die Berechnung der Schatten das sogenannte \emph{Shadow Mapping} verwendet. Dazu wird eine zweite Kamera implementiert, die die Spielwelt aus der Sicht einer Lichtquelle, im Fall der Blocklib die Sonne, betrachtet. Die so gerenderte Sicht wird dann genutzt, um zu berechnen ob ein Fragment das erste Hindernis aus Sicht der Lichtquelle ist. Ist dies der Fall, so wird das Fragment beleuchtet, ansonsten ist etwas anderes vor dem Fragment und es befindet sich im Schatten.

Wird die \class{ShadowMap} mitangezeigt, so kann man erkennen, dass diese sich zeitgleich mit dem Flackern verschiebt. Die Position der \class{ShadowMap} wird über die Klasse \class{ShadowBounds} bestimmt. Die Position muss in zwei Fällen geändert werden, zum einen, wenn die Sonne sich bewegt, also, wenn der Tag-Nacht-Zyklus aktiv ist, zum anderen, wenn sich die Spielerkamera bewegt, da die Darstellung von Schatten immer mit dem Sichtfeld des Spielers übereinstimmen muss.

Die Ausgabe der Position der \class{ShadowBounds} während des Spiels bei deaktiviertem Tag-Nacht-Zyklus und stillstehender Spielerkamera bestätigt, dass diese Position sich ztatsächlich zeitgleich mit dem Auftreten des Flackerns verändert. Das Auftreten der Wettkampfbedingung lässt sich also erkennen, indem eine Änderung der Position der \class{ShadowBounds} in aufeinanderfolgenden Frames erkannt wird.

Nun gilt es die Ursache der Positionsänderung ausfindig zu machen. Die Position der \class{ShadowBounds} hängt selbst von mehreren Variablen ab. Verfolgt man den Verlauf der Änderungen, finden sich die folgenden Abhängigkeiten:

\begin{tabular}{ll}
	\class{ShadowMap} &$\to$ \code{ShadowBounds.update}\\
	& $\to$ \class{LightViewMatrix}\\
	& $\to$ \code{DayNightLighting.getSunUp()} \\
	& $\to$ \code{DayNightLighting.position}\\
	& $\to$ \code{DayNightLightig.updateLightPosition(float, boolean)}
\end{tabular}

\code{DayNightLightig.updateLightPosition(float, boolean)} wird nicht im Renderthread ausgeführt sondern während \code{update()} in einem anderen Thread. Der relevante Abschnitt der Methode ist 
\begin{lstlisting}[]
private void updateLightPosition(float progress, boolean day) {
	// ...
	position = new Vector3f(direction.x, direction.y, direction.z);
	position.scale(SUN_HEIGHT); (*\label{lst:updateLightPosition:scale}*)
	// let sun stay relative to the player
	position = Vector3f.add(position, Context.getInstance().getCamera().getPosition(), new Vector3f());
}
\end{lstlisting}
Wird also der Wert von \code{DayNightLighting.position} im Renderthread ausgelesen, während die Methode in Ausführung ist, beispielsweise während Zeile~\ref{lst:updateLightPosition:scale}, so enthält \code{DayNightLighting.position} einen vollkommen falschen Wert. Hier existiert also die gesuchte Wettkampfbedingung.
\subsection{Umstellung auf zustandsloses Rendering}
In der Blocklib implementieren alle zu zeichnenden Elemente das Interface \class{Renderable}. Dieses definiert die Funktionalität, die zum Zeichnen eines Elements notwendig ist, sowie die Methoden \code{show()} und \code{hide()}. Das Rendersystem ist so aufgebaut, dass ein Element nach einem Aufruf von \code{show()} solange gezeichnet wird, bis \code{hide()} aufgerufen wird. Dies wird erreicht, indem eine Datenstruktur im \class{MasterRenderer} alle Renderables speichert. Diese zustandsbehaftete Zeichenmethode birgt Vor- und Nachteile.

\todo{itemize mit + und - ?}
\begin{itemize}
	\item[$+$] Da \code{show()} nur einmal aufgerufen werden muss, können auch Systeme ohne Updatemethode das Zeichnen von Elementen veranlassen.
	\item[$+$] Da die Elemente gespeichert sind, müssen sie nicht jedes Mal neu hinzugefügt werden. Das verringert den Rechenbedarf.
	\item[$-$] Werden Renderables häufig ausgetauscht, müssen die alten Elemente jedes Mal entfernt werden.
	\item[$-$] Da die Einführung eines parallelen Renderthreads ansonsten zu Race Conditions führen würde, muss diese Datenstruktur aus Sicht des Renderthreads während des Zeichnens unverändert sein. Wie beschrieben, wird dazu ein Double Buffer eingesetzt. Bei einem zustandsbehafteten Double Buffer müssen beim Swap die Elemente des einen Buffers tatsächlich in den anderen Buffer kopiert werden. Das erfordert Zeit, die nicht parallelisiert werden kann, da das Wechseln des Buffers synchronisiert sein muss.
	\item[$-$] Die Existenz von paarweise aufzurufenden Funktionen birgt die Gefahr, dass der zweite Aufruf vergessen wird. Wie bei \code{malloc()} und \code{free()} in C entsteht durch einen fehlenden Aufruf von \code{hide()} ein Speicherleck. Zudem wird dann die Anzahl der zu zeichnenden Elemente immer größer und der Rendervorgang wird verlangsamt. Des weiteren werden möglicherweise Renderables gezeichnet, die nicht gezeichnet werden sollen.
\end{itemize}

Um den Zeichenaufwand zu veringern, implementiert die Klasse \class{ChunkManager} das sogenannte \emph{Frustum Culling}. Es wird berechnet, welche Chunks sich im Sichtfeld der Kamera befinden. Nur diese Chunks sollen gezeichnet werden. Dazu entfernt der ChunkManager jeden Frame alle Chunks mittels \code{hide()} aus der Datenstruktur der zu zeichnenden Elemente und fügt nur die als sichtbar ermittelten Chunks wieder ein. Die Chunks machen mit etwa 
% 334*2/(334*2+4+73+142)
75 \% bis
% 529*2/(529*2+4+73+142)
82 \% einen Großteil aller Rendereables aus. Es wird also bereits ein Großteil der zu zeichnenden Elemente in jedem Frame neu hinzugefügt. Des weiteren gibt es in der Blocklib bis jetzt keine Klasse, die Renderables zu der Datenstruktur hinzufügt, aber keine Updatemethode besitzt. Somit werden die Gefahr von vergessenen \code{hide()} Aufrufen und der Kopieraufwand des Double Buffers als wichtiger eingeschätzt, als die oben beschriebenen Vorteile.

Die Datenstruktur ist nun wie folgt implementiert. Die Methode \code{hide()} entfällt ersatzlos. Stattdessen wird nach dem Swap des Double Buffers der Renderables, der Buffer für die nächsten zu zeichnenden Elemente geleert, sodass Renderables automatisch nicht mehr gezeichnet werden, wenn sie nicht mehr hinzugefügt werden. Die Methode \code{show()} wird zu \code{draw()} umbenannt, um zu signalisieren, dass es sich um einen einmaligen Vorgang handelt. Alle bisherigen Aufrufe von \code{show()} werden durch \code{draw()} in Updatemethoden ersetzt, sodass das Verhalten zum vorherigen Stand identisch ist.
\section{Integration und Performanceanalyse}

\chapter{Fazit und Ausblick}
\section{Etablierung einer Composition Root}
\begin{figure}
  \begin{center}
    \includesvg[]{Context-shortened.svg}
  \end{center}
  \caption{A}\label{fig:diagContext}
\end{figure}
\textcite[S.~76]{Seemann2012} definiert die \gls{compositionRoot} als (vorzugsweise) genau eine Position in einer Anwendung an der Module zusammengesetzt werden. Die korrekte Nutzung einer \gls{compositionRoot} impliziert die weite Verbreitung von \gls{dependencyInjection}, der Designarchitektur, in der Module ihre Abhängigkeiten nicht selbst erzeugen, sondern diese Aufgabe an aufrufende Module übertragen. Schlussendlich wird diese Aufgabe bis zur \gls{compositionRoot} weitergegeben, wo alle Abhängigkeiten aufgelöst werden. 


Das Design vieler Module in der Blocklib verfolgt häufig einen anderen Ansatz. Insbesondere Die abstrakte Klasse \class{Context} und die von ihr erbenden Klassen entsprechen eher einer Mischung aus \gls{servicelocator}~\cite[S.~301~ff.]{Nystrom2015} und \gls{singleton}~\cite[S.~103~ff.]{Nystrom2015}. Ein Klassendiagramm von \class{Context} ist in Abbildung~\vref{fig:diagContext} dargestellt. Die statische Methode \code{getInstance()} weist auf das \gls{singleton} Pattern hin. Die Klasse besitzt zwar keine Methode \code{resolve(...)}, wie es bei einem üblichen \gls{servicelocator} der Fall wäre, bietet aber stattdessen konkrete Methoden für die einzelnen Dienste, die es in der Blocklib gibt.

Sowohl das ursprüngliche \gls{singleton} Pattern von~\textcite[S.~127~ff.]{Gamma2016}, als auch das \gls{servicelocator} Pattern werden inzwischen teils kritisch betrachtet~\cite[S.~103~ff.]{Nystrom2015}\cite[S.~154~ff.]{Seemann2012}. Das liegt insbesondere daran, dass es beide Patterns ermöglichen, Module über die gesamte Code Basis hinweg zu koppeln, was es wiederum erschwert Auswirkungen von Änderungen im Code abzuschätzen. Des Weiteren erschwert besonders das \gls{singleton} Pattern von \class{Context} in der Blocklib das (Unit-) Testen. Nutzt eine Klasse \class{Context}, so ist es nahezu unmöglich, diese isoliert zu testen, da deren Methodenaufrufe sich durch den \class{Context} über die ganze Code Bases erstrecken können. Diese  Aufrufe können auch nicht durch \glspl{testdouble} (also spezielle Module, die das Verhalten von Komponenten vorspielen, damit die Originale beim Testen nicht benötigt werden) ersetzt werden, weil \class{Context} keine Möglichkeit bietet die bereitgestellten Dienste zu ersetzen. Außerdem kann im Test nicht festgestellt werden, welche Methoden überhaupt ersetzt werden müssten, da die Abhängigkeiten einer Klasse versteckt sind. Ruft eine Methode beispielsweise intern \code{Context.getInstance().getWorldInteraction()} auf, so ist das von außen nicht zu erkennen.

\subsubsection{Vorteile einer \glsentrytext{compositionRoot}}

Die Nutzung eines \gls{singleton} und \gls{servicelocator} ist zwar anfangs sehr angenehm, da jedes Modul nur einen Methodenaufruf entfernt ist, sie bringt aber auch die gleichen Probleme mit sich wie globale Variablen. Anders verhält es sich bei der Nutzung von \gls{dependencyInjection} über Konstruktoren und einer \gls{compositionRoot}. Im Folgenden sind kurz einige ihrer Vorteile erläutert.

\paragraph{Erkennung von Abhängigkeiten}Abhängigkeiten werden sichtbar, wenn die Klasse einen Konstruktor besitzt, der eine Instanz von \class{WorldInteraction} erwartet. Dann ist die Abhängigkeit offensichtlich. Es kann also einfach abgeschätzt werden, welche Änderungen Einfluss auf ein bestimmtes Modul haben. Zudem ist klar welche Bereiche beeinflusst werden können, wenn Methoden der Klasse aufgerufen werden, nämlich nur diejenigen zu denen sie auch eine Referenz besitzt.

\paragraph{Keine versehentlichen Abhängigkeiten} Außerdem wird die Gefahr verhindert aus Versehen, neue Abhängigkeiten aufzubauen, denn diese müssen explizit im Konstruktor angegeben werden. Bei einem \gls{servicelocator} reicht ein einfacher Methodenaufruf.

\paragraph{Das Testen wird vereinfacht} Es ist klar, dass die Klasse \class{WorldInteraction} genutzt wird, sie wird ja schließlich dem Konstruktor übergeben. Der Austausch von \class{WorldInteraction} durch ein Test Double ist trivial, da dieses einfach beim Aufruf des Konstruktors der Klasse übergeben werden kann. Dadurch lässt sich einfach eines der großen Probleme der Blocklib angehen -- eine fehlende automatisierte Testabdeckung. Aktuell muss bei den meisten Tests die gesamte Blocklib gestartet werden, sodass die Tests zu lange dauern, als dass man sie in häufigen Intervallen durchführen kann.

\textcite[S.~15~ff.]{Seemann2012} beschreibt noch einige weitere Vorteile der Nutzung von \gls{dependencyInjection}, darunter Erweiterbarkeit und die Möglichkeit leichter parallel Code zu entwickeln. Ein besonders wichtiger Punkt für die Blocklib ist zudem die verbesserte Wartbarkeit, da der Durchsatz an Entwicklerinnen und Entwicklern, die an der Blocklib arbeiten, sehr hoch ist.

\subsubsection{Mögliche Probleme bei der Umstellung}

Somit kann man davon ausgehen, dass die breite Etablierung der Nutzung von \gls{dependencyInjection} in Verbindung mit einer zentralen \gls{compositionRoot} erstrebenswert ist. Das Wort \enquote{breit} lässt allerdings bereits vermuten, dass diese Etablierung sehr viele Änderungen in der Blocklib hervorrufen wird. Im Folgenden wird nun beschrieben, welche potentiellen Probleme dabei in der Blocklib zu bewältigen sind.

\paragraph{Wo ist die \glsentrytext{compositionRoot}?}
Typischerweise soll die Komposition des Objektgraphen, also die Zusammensetzung der Module in der \gls{compositionRoot}, so nah wie möglich am Einstiegspunkt des Programms liegen~\cite[S.~232~ff.]{Martin17}\cite[S.~76~f.]{Seemann2012}. Das liegt daran, dass man dadurch die Entscheidung, wie die Komposition auszusehen hat und zugleich konkrete Abhängigkeiten, so lange wie möglich herauszögern kann. Man erreicht so also eine möglichst lose Kopplung der einzelnen Module des Programms.

Die Blocklib stellt nun einen besonderen Fall dar, da sie eigentlich nur eine Bibliothek darstellt. Der Programmeinstiegspunkt liegt also an einer anderen Stelle, nämlich in dem \class{BlockLibProgram}, das ausgeführt wird. Platziert man die \gls{compositionRoot} dort, besitzt jedes Programm völlige Freiheit, bei der Komposition der Blocklib. Allerdings ist dies dem Ziel abträglich, die Blocklib möglichst einfach nutzen zu können, zumal in den allermeisten Fällen die Komposition identisch sein wird. Einen Sonderfall stellt hier die Nutzung der Blocklib in der Multiplayer-Version dar, da hier tatsächlich beispielsweise für die Klassen \class{WorldInteraction} und \class{Context} spezielle Unterklassen vorhanden sind, die je nach Blocklib-Version verwendet werden. Alternativ kann man die Klasse \class{Game} als Einstiegspunkt betrachten, da hier bereit einige bibliotheksinterne Module komponiert werden und sie die Klasse ist, die in dem \class{BlockLibProgram} instanziiert wird. Sie stellt also den Einstiegspunkt für die Bibliothek dar. Wird die \gls{compositionRoot} allerdings hier platziert, verliert das \class{BlockLibProgram} allerdings die Möglichkeit der freien Konfiguration.

Wird die Blocklib in der Multiplayer-Version genutzt, entsteht ein weiteres Problem, da der gestartete Spring-Server keine Parameter übergeben bekommt. Es gibt nur einen Standardkonstruktor. Es muss also evaluiert werden, ob und wie Abhängigkeiten an den Server übergeben werden können. Ansonsten ist er von der \gls{compositionRoot} ausgeschlossen.

\paragraph{Aufwand der Umstellung ist nicht abzuschätzen}
Wie bereits erwähnt, wird die Etablierung einer \gls{compositionRoot} in der Blocklib sehr viele Änderungen hervorrufen, da beinahe jedes Modul über den \class{Context} verfügbar ist. Es lässt sich allerdings nicht abschätzen, wie hoch der Aufwand für eine Umstellung tatsächlich ist. Es gibt vielleicht Stellen, an denen das Design eines Moduls davon abhängt, dass es einen \class{Context} gibt, sodass das Modul komplett umstrukturiert werden muss. Andererseits kann es auch sein, dass die Umstellung nur die Erzeugung von geeigneten Konstruktoren erfordert und sehr simpel ist. Aufgrund der durch den \class{Context} versteckten Abhängigkeiten lässt sich das im Voraus nicht ermitteln, sondern wird sich erst bei der Durchführung zeigen.

\section{Dynamisches Rendern von Schrift}
\section{Thread Affinity}
\printnoidxglossaries

\printbibliography[title={Literaturverzeichnis},heading=bibintoc,notkeyword=online]

\printbibliography[title={Quellenverzeichnis},heading=bibintoc,keyword=online] 


\end{document}